\RequirePackage{fix-cm}
\documentclass[a4paper,11pt]{ltjsarticle}
\usepackage{amsmath,amssymb,amsthm,mathtools,bm}
\usepackage[margin=25mm]{geometry}
\usepackage{booktabs,longtable,array}
\usepackage[unicode,colorlinks=true,linkcolor=blue,citecolor=blue,urlcolor=blue]{hyperref}
\usepackage{luatexja-fontspec}
\setmainjfont{Noto Serif CJK JP}
\setsansjfont{Noto Sans CJK JP}
\numberwithin{equation}{section}
\theoremstyle{plain}
\newtheorem{theorem}{定理}[section]
\newtheorem{proposition}[theorem]{命題}
\newtheorem{lemma}[theorem]{補題}
\newtheorem{corollary}[theorem]{系}
\theoremstyle{definition}
\newtheorem{definition}[theorem]{定義}
\newtheorem{remark}[theorem]{注意}
\renewcommand{\proofname}{証明}
\newcommand{\C}{\mathbb C}
\newcommand{\A}{\mathbb A}
\newcommand{\Pj}{\mathbb P}
\newcommand{\Gr}{\operatorname{Gr}}
\newcommand{\Pois}{\operatorname{Pois}}
\newcommand{\dd}{\mathrm d}
\newcommand{\Version}{1.2}
\title{次数1・2の del Pezzo manifold の\texorpdfstring{\\}{ }Poisson 構造とモジュライ}
\author{}
\date{2026年11月20日\quad v\Version\quad 公表用原稿}
\hypersetup{
 pdftitle={次数1・2の del Pezzo manifold の Poisson 構造とモジュライ},
 pdfauthor={},
 pdfsubject={Poisson schemes, relative classification, and moduli; v1.1}
}
\begin{document}
\maketitle
\begin{abstract}
滑らかな複素 $n$ 次元 del Pezzo manifold $X$、$n\geq2$、
$-K_X=(n-1)H$、$H^n\in\{1,2\}$ の大域 Poisson 構造を扱う。
重み付き Euler 列と偏微分列の Koszul 完全性から全ての二重ベクトルを記述し、
Jacobi 方程式の零スキームが次数1では $\mathbb A^n$、次数2では
$\operatorname{Gr}(2,n+1)$ の Pl\"ucker affine cone に等しいことを証明する。
この同定を滑らかな del Pezzo manifold の族へ拡張し、
非被約基底を含む任意の基底変換との両立を証明する。
座標変換に伴う線束の twist を明示し、Poisson 構造付きの族を分類する
有限型 Deligne--Mumford スタックを構成する。
次数2の零切断の blow-up によって滑らかなモジュライスタックを得て、
Pl\"ucker 錐の特異性を次元別に分類する。
Poisson 構造を忘れる射は affine かつ平坦であり、
次数1では階数 $n$ のベクトル束、次数2では相対 Pl\"ucker cone である。
非零構造の一般階数は2であり、射影化は
$\operatorname{Gr}(n-1,H^0(X,H))$ に内在的に同定される。
標準 ambient stack への一意な Poisson lift、葉の記述、曲面での退化因子、
中心元による次数付き形式量子化を併せて示す。
\end{abstract}
\tableofcontents

\section{問題と主結果}
\label{sec:introduction}
本稿の対象は、滑らかな連結複素射影多様体 $X$ と ample Cartier 因子 $H$ で
\[
 n=\dim X\geq2,\qquad -K_X\simeq(n-1)H,\qquad d=H^n\in\{1,2\}
\]
を満たすものである。固定した $X$ 上の分類から始め、
$X$ 自体も変化する族とそのモジュライを構成する。
$V_X=H^0(X,\bigwedge^2T_X)$、$W_X=H^0(X,\bigwedge^3T_X)$ とおき、
二次写像 $Q_X:V_X\to W_X$、$Q_X(\pi)=\frac12[\pi,\pi]$ の零スキームを
\[
 \Pois_{\mathrm{sch}}(X)=
 \operatorname{Spec}\!\left(\operatorname{Sym}(V_X^*)/
       (Q_X^*(W_X^*))\right)
\]
と書く。ここで括弧は Schouten 括弧であり、右辺のイデアルは
$Q_X$ の全ての係数で生成する。自己同型群による商は取らない。

\begin{theorem}[固定多様体上の Poisson スキーム]
\label{thm:main}
上の仮定の下で、次が成り立つ。
\begin{enumerate}
\item $d=1$ なら $\dim V_X=n$ であり、$Q_X$ は恒等的に零である。
従って $\Pois_{\mathrm{sch}}(X)\simeq\mathbb A^n$。
\item $d=2$ なら $\dim V_X=\binom{n+1}{2}$ であり、適切な線型座標の下で
Jacobi イデアルは Pl\"ucker イデアルに等しい。
従って $\Pois_{\mathrm{sch}}(X)$ は
$\Gr(2,n+1)\subset\mathbb P^{\binom{n+1}{2}-1}$ の affine cone である。
これらは被約かつ既約なスキームである。
\item 全ての非零複素 Poisson 構造は一般点で階数2を持ち、スカラーによる射影化は
\[
 \mathbb P\Pois(X)\simeq\Gr(n-1,H^0(X,H))
\]
に内在的に同定される。
\item 標準埋め込み $X\subset\mathcal P$、
$\mathcal P=\mathbb P(1^n,2,3)$ または $\mathbb P(1^{n+1},2)$ を
重み付き射影スタックとして解釈すると、各 Poisson 構造は
$X$ を Poisson 部分多様体とする $\mathcal P$ 上の Poisson 構造に一意に延長する。
定義式 $F$ を固定すれば、cone の多項式環上で $F$ を Casimir とする
重み0の代表も一意である。
\end{enumerate}
\end{theorem}

ここで $\mathbb A^n$ や affine cone との線型同定には座標の選択が入る。
一方、射影 Grassmann 多様体による記述は、$K_X^{-1}\simeq H^{\otimes(n-1)}$
の定数倍の選択を除けば内在的である。
スキームとしての主張は、任意の可換 $\C$ 代数 $B$ に対し、
$V_X\otimes B$ 内で Jacobi 恒等式を満たす元が同じ方程式で分類されることを意味する。
動く $X$ に対する主張には、さらに基底変換と座標変換に関する議論が必要となる。

\begin{theorem}[族と Poisson モジュライ]
\label{thm:relative-main}
$n\geq2$ と $d\in\{1,2\}$ を固定する。
滑らかな $n$ 次元・次数 $d$ の del Pezzo manifold の
モジュライを $\mathcal D_{n,d}$ とする。
相対 Poisson 二重ベクトルを加えたモジュライ $\mathcal P_{n,d}$ は
有限型 Deligne--Mumford スタックであり、忘却射
\[
 \mathcal P_{n,d}\longrightarrow\mathcal D_{n,d}
\]
は affine、有限表示かつ平坦である。
次数1では階数 $n$ のベクトル束、次数2では局所的に
$\operatorname{Gr}(2,n+1)$ の Pl\"ucker affine cone の族となり、
この記述は任意の基底変換と両立する。
次数1の全体、および次数2の非零 Poisson 構造の部分は滑らかである。
次数2では $n=2$ のとき全体が滑らかで、$n\geq3$ のとき特異点の集合は零切断に等しい。
非零構造をスカラー倍で割った相対モジュライは
基本線束を局所的に選んで得る相対 Grassmann 多様体の降下である。
\end{theorem}

ここで基本偏極は線束そのものに付随するスカラー自己同型を除いて扱う。
基本偏極のクラスは族に一意に定まり、追加の標識を必要としない
（補題~\ref{lemma:fundamental-polarization-class}）。
具体的な対象、射、および商表示は定理~\ref{thm:poisson-moduli} で定義する。
特に、$\mathcal D_{n,d}$ 上の大域的な自明直積や基本線束の大域的存在は主張していない。
族 $p:\mathcal X\to S$ 上で実際の基本線束 $H$ を選んだときの大域二重ベクトル束は
\[
 p_*\bigwedge^2T_{\mathcal X/S}
 \simeq\bigwedge^{n-1}(p_*H)\otimes
 p_*(\omega_{\mathcal X/S}^{-1}\otimes H^{-(n-1)})
\]
である。右辺の最後の線束が、同変性に必要な twist を担う。
この同型とその降下は命題~\ref{prop:intrinsic-bivector-bundle} で証明する。

\subsection{先行研究との関係と到達範囲}
標準重み付きモデルは del Pezzo 多様体の既知の分類による。
$n\geq3$ では Kuznetsov--Prokhorov \cite[Theorem 1.2(i),(ii)]{KP2023}、
曲面では Cheltsov--Prokhorov \cite[Theorem 2.5]{CP2021} の滑らかな場合を用いる。
以下の大域二重ベクトルの計算は、これらのモデルを入力として直接行う。

3次元 Fano 多様体上の Poisson 構造は Loray--Pereira--Touzet
\cite[Theorems 7.1, 9.1]{LPT2013} により研究されている。
次数1・2の del Pezzo threefold に現れる成分は基本線束 $H$ の pencil に由来し、
本稿の射影パラメータ空間はそれぞれ $\mathbb P^2$、$\Gr(2,4)$ である。
この3次元の集合論的分類を本稿の新しい結果として扱わない。
Mayanskiy \cite[Introduction, Theorem 1]{Mayanskiy2015} は cubic threefold について
Pl\"ucker 関係と Poisson スキームを明示的に計算している。
従って本稿で比較すべき点は、次数1・2、任意次元での大域計算と、
次数2における Jacobi イデアル自体の回収である。
これらの主張の優先権を確定する文献調査は継続中である。
また曲面について Hong--Xu \cite[Theorem 3.9]{HongXu2011} は
一般位置の $r$ 点で吹き上げた $\mathbb P^2$、$5\leq r\leq8$ の
Poisson cohomology を計算している。
次数1・2の $\dim H^2_\pi=10,9$ および $H^1_\pi=0$ は既知であり、
本稿の曲面モジュライの次元とも整合する。
本稿の族に関する結果は、これらの数値の新しい発見ではなく、
任意次元での Jacobi 零スキーム、非被約基底を含む基底変換、
線束の twist と具体的商表示を一つの構成として与えるものである。
有限型商スタックおよび基底変換については
Stacks Project \cite[Tags 0B91, 06FI, 0DSL]{StacksProject} の一般結果を用いる。

形式量子化の存在には Kontsevich の affine formality \cite{Kontsevich} を用いる。
曲面での cyclic potential と中心元は
Etingof--Ginzburg \cite[Lemma 7.4.2, Proposition 8.1.1(i)]{EG2010} に基づく。
本稿の中心的超曲面表示はこれら既知の定理の適用であり、
量子化の一意性や非可換 del Pezzo 多様体全体の分類を意味しない。

\subsection{構成}
del Pezzo manifold の定義と標準模型を短く整理した後、
全ての二重ベクトルを Casimir 代表に直して Jacobi イデアルを決定する。
続いて内在的な Grassmann 表示、symplectic leaf、曲面の反標準零因子を調べる。
これを相対化して Poisson モジュライを構成し、
その零切断の blow-up と特異性を調べ、その後に形式量子化と曲面の具体表示を扱う。
量子化の代数化および条件付き AS 正則性は、可換な Poisson モジュライの構成とは区別する。
次数3・4の詳細な分類は本稿に含めず、関連する構成と障害を研究ノートに置く。

% !TeX root = ../poisson-del-pezzo.tex
\section{del Pezzo manifold の概観}
\label{sec:del-pezzo-review}

\subsection{基本偏極と次数}
本稿で del Pezzo manifold とは，滑らかな連結複素射影多様体 $X$ と
ample 線束 $H$ で
\[
 n=\dim X\geq2,\qquad \omega_X^{-1}\simeq H^{\otimes(n-1)}
\]
を満たすものをいう．$H$ を基本偏極，$d=\int_Xc_1(H)^n$ を次数と呼ぶ．
これは Fujita の定義に従うものであり
\cite[Theorem (1.9), Definition (1.13)]{Fujita1980}，
$n=2$ では通常の del Pezzo 曲面と $H=\omega_X^{-1}$ に一致する．
$H$ が原始的なら指数は $n-1$，余指数は $2$ である．
ここでは原始性を定義に含めないので，例えば
$(\Pj^3,\mathcal O_{\Pj^3}(2))$ も次数 $8$ の対象に入る．
Picard 数 $1$ という条件も課さない．

基本線型系には
\begin{equation}
 h^0(X,H)=n+d-1
 \label{eq:fundamental-sections}
\end{equation}
が成り立つ．次数 $1$ では $|H|$ の基底スキームは一つの被約点であり，
次数 $2$ では $|H|$ は分岐因子が滑らかな四次超曲面である
二重被覆 $X\to\Pj^n$ を定める．次数 $d\geq3$ では $H$ は非常に ample である
\cite[Lemma (1.10), Theorem (2.5)]{Fujita1980}
\cite[Proposition 2.2]{KP2023}．
また一般の基本因子による切断を繰り返すと，同じ次数の del Pezzo 曲面に至る．
さらに $n-1$ 個の一般の基本因子の完全交叉 $C$ は滑らかな曲線であり，
随伴公式から
$\omega_C\simeq(\omega_X\otimes H^{\otimes(n-1)})|_C\simeq\mathcal O_C$
を得る．この種数 $1$ の切断が，曲面と高次元の場合を結び付ける．

\subsection{標準模型}
滑らかな場合の分類は Fujita による．
低次数の超曲面表示については
\cite[Theorems (2.1), (2.2), (2.5)]{Fujita1980} と
\cite[Theorems (14.6), (16.7)]{Fujita1984} を，
高次元の一覧については \cite[Theorem 1.2]{KP2023} を参照する．
$n\geq3$ の場合の模型は次の通りである．表の線型切断は期待次元を持つ滑らかな切断とし，
記載した偏極を用いる．
\begin{center}
\begin{tabular}{@{}cll@{}}
\toprule
$d$ & $X$ と基本偏極 $H$ & 次元の範囲\\
\midrule
$1$ & $X_6\subset\Pj(1^n,2,3)$，$\mathcal O_X(1)$ & $n\geq3$\\
$2$ & $X_4\subset\Pj(1^{n+1},2)$，$\mathcal O_X(1)$ & $n\geq3$\\
$3$ & $X_3\subset\Pj^{n+1}$，$\mathcal O_X(1)$ & $n\geq3$\\
$4$ & $X_{2,2}\subset\Pj^{n+2}$，$\mathcal O_X(1)$ & $n\geq3$\\
$5$ & $\Gr(2,5)\cap\Pj^{n+3}\subset\Pj^9$，$\mathcal O_X(1)$ & $3\leq n\leq6$\\
$6$ & $(\Pj^2\times\Pj^2)\cap\Pj^{n+4}\subset\Pj^8$，$\mathcal O_X(1,1)$ & $n=3,4$\\
 & $(\Pj^1)^3$，$\mathcal O(1,1,1)$ & $n=3$\\
$7$ & $\operatorname{Bl}_{p}\Pj^3$，$2L-E$ & $n=3$\\
$8$ & $\Pj^3$，$\mathcal O(2)$ & $n=3$\\
\bottomrule
\end{tabular}
\end{center}
ここで $L$ は超平面束の引き戻し，$E$ は例外因子である．
最初の二つの模型は ambient の非自明な stabilizer を持つ点を避ける．
曲面では $1\leq d\leq9$ であり，$\Pj^2$，$\Pj^1\times\Pj^1$，
および $\Pj^2$ の一般位置にある高々 $8$ 点の blow-up が現れる
\cite[\S3]{Fujita1980}．
次数 $1,2$ の曲面にも表と同じ重み付き模型が成り立つ
\cite[Theorem 2.5(v)--(vii)]{CP2021}．

本稿で用いる基本切断環は
$R(X,H)=\bigoplus_{r\geq0}H^0(X,H^{\otimes r})$ である．
次数 $1$ では重み $(1^n,2,3)$ の生成元と重み $6$ の一つの関係，
次数 $2$ では重み $(1^{n+1},2)$ の生成元と重み $4$ の一つの関係を持つ．
従って Hilbert 級数はそれぞれ
\begin{equation}
 \frac{1-t^6}{(1-t)^n(1-t^2)(1-t^3)},
 \qquad
 \frac{1-t^4}{(1-t)^{n+1}(1-t^2)}
 \label{eq:del-pezzo-hilbert-series}
\end{equation}
である．これらは次数付き超曲面の完全列から直接従う．
族を扱う際には，この切断環の相対版が座標付き方程式と偏極多様体をつなぐ．
ただし，各ファイバーで基本偏極の同型類が一意であることと，
族全体で線束および $\omega^{-1}\simeq H^{\otimes(n-1)}$ の同型を選ぶことは異なる．
この選択に伴う基底上の線束とスカラー自己同型は，モジュライ構成で明示する．

\subsection{Poisson 問題の位置付け}
曲面では $\bigwedge^2T_X\simeq\omega_X^{-1}$ かつ
$\bigwedge^3T_X=0$ なので，全ての反標準切断が Poisson 構造を定める．
従って次数 $1,2$ の固定曲面上では，それぞれ $\A^2,\A^3$ が得られる．
曲面の Poisson cohomology は Hong--Xu \cite{HongXu2011} により計算されている．
特に同論文の Theorem 3.9 では，次数 $1,2$ の場合の
$H^1_\pi(X)=0$，$\dim H^2_\pi(X)=10,9$，$H^i_\pi(X)=0$（$i>2$）が示される．
これは多様体も動かす無限小変形についての既知の比較対象である．
三次元での分類は \cite[Theorems 7.1, 9.1]{LPT2013} に，
cubic threefold の Jacobi 方程式の明示的計算は
\cite[Theorem 1 and the subsequent Remark]{Mayanskiy2015} にある．
以下ではこれらを踏まえ，次数 $1,2$ の任意次元での
スキーム論的分類と，基底変換・座標変更を伴う族への拡張を扱う．
次数 $3$ 以上の Poisson 構造を同じ模型で分類することは本稿の主張に含めない．

% !TeX root = ../poisson-del-pezzo.tex
\section{重み付きモデルと大域的多重ベクトル}
\label{sec:models}

以下，基礎体は $\mathbb C$ とする．$X$ は滑らかな連結射影多様体，
$\dim X=n\geq 2$，$H$ は ample Cartier 因子であり，
$-K_X\simeq (n-1)H$ および $d=H^n\in\{1,2\}$ を満たすとする．
ここでのすべての Poisson 構造は代数的なものを指す．
$n\geq3$ では del Pezzo 多様体の重み付きモデル
\cite[Theorem 1.2]{KP2023}，$n=2$ では反標準モデル
\cite[Theorem 2.5]{CP2021} により，次の方程式を取れる．
\begin{align}
 d=1:\quad
 S&=\mathbb C[x_0,\ldots,x_{n-1},y,z],
 &F&=z^2-y^3-A_4(x)y-A_6(x),
 &\deg(x_i,y,z)&=(1,2,3),
 \label{eq:degree-one-model}\\
 d=2:\quad
 S&=\mathbb C[x_0,\ldots,x_n,z],
 &F&=z^2-f_4(x),
 &\deg(x_i,z)&=(1,2).
 \label{eq:degree-two-model}
\end{align}
添字は多項式の次数を表す．$F$ の重み付き次数を $m$ と書くと，
それぞれ $m=6,4$ である．平方完成と $y$ の平行移動は標数 $0$ で許される．
$R=S/(F)$ は $H$ に関する全切断環であり，$H=\mathcal O_X(1)$ である．

変数を一括して $q_0,\ldots,q_{N-1}$，その重みを $w_i$ と書く．
ここで $N=n+2$ である．ambient は重み付き射影\emph{スタック}
\[
 \mathcal P=[(\mathbb A^N\setminus\{0\})/\mathbb G_m]
\]
とする．次数 $1$ の非自明な stabilizer は $y$ 軸と $z$ 軸に，
次数 $2$ のそれは $z$ 軸にのみ現れ，\eqref{eq:degree-one-model}，
\eqref{eq:degree-two-model} の $X$ はこれらを避ける．
したがって $\mathcal P$ 内の $X$ は元の滑らかなスキームである．
次数 $2$ では $X$ の滑らかさから分岐四次超曲面
$(f_4=0)\subset\mathbb P^n$ も滑らかである．

$\partial_{q_i}$ に重み $-w_i$ を与え，
$\operatorname{Der}S=\operatorname{Der}_{\mathbb C}S$ と略記する．
Euler ベクトル場と縮約の規約を
\[
 E=\sum_i w_iq_i\partial_{q_i},\qquad
 \iota_\alpha(v_1\wedge\cdots\wedge v_r)
 =\sum_{j=1}^r(-1)^{j-1}\alpha(v_j)
 v_1\wedge\cdots\wedge\widehat v_j\wedge\cdots\wedge v_r
\]
で定める．重み $0$ の二重ベクトル $B$ は $\mathcal P$ 上に降下し，
$B$ と $B+E\wedge V$ は同じ切断を表す．ここで $V\in(\operatorname{Der}S)_0$ である．
括弧は $\{a,b\}=B(\mathrm da,\mathrm db)$ とする．
$F$ が Casimir である条件 $\iota_{\mathrm dF}B=0$ は，
単に $(F)$ が Poisson ideal となる条件より強い．

\begin{lemma}[必要な消滅]
\label{lemma:weighted-vanishing}
任意の整数 $t$ に対して
\[
 H^0(\mathcal P,\mathcal O_{\mathcal P}(t))=S_t,
 \qquad
 H^i(\mathcal P,\mathcal O_{\mathcal P}(t))=0
 \quad(0<i<N-1),
\]
および
\[
 H^{N-1}(\mathcal P,\mathcal O_{\mathcal P}(t))
 \simeq S_{-t-\sum_iw_i}^{\vee}
\]
が成り立つ．したがって
\[
 H^0(X,\mathcal O_X(t))=R_t,\qquad
 H^1(X,\mathcal O_X(t))=0,\qquad H^2(X,\mathcal O_X)=0.
\]
\end{lemma}

\begin{proof}
開集合 $q_i\ne0$ による \v Cech 複体を Laurent 単項式ごとに分解する．
負指数を持つ変数の集合が空でも全体でもない場合，対応する部分複体は完全である．
負指数を持つ変数がない場合が次数 $0$ に，すべての変数の指数が負の場合が
次数 $N-1$ に寄与する．後者の単項式を $(q_0\cdots q_{N-1})^{-1}$ と比較すれば
表示した双対性を得る．$\mathbb G_m$ の重み成分を取る操作は完全である．

超曲面の完全列
\[
 0\longrightarrow\mathcal O_{\mathcal P}(t-m)
 \xrightarrow{\ F\ }\mathcal O_{\mathcal P}(t)
 \longrightarrow\mathcal O_X(t)\longrightarrow0
\]
と $N-1=n+1\geq3$ から $H^0$ の記述と $H^1$ の消滅を得る．
$n\geq3$ では同じ列から $H^2(X,\mathcal O_X)=0$ も従う．
$n=2$ で問題となる ambient の最上コホモロジーは
$S_{m-\sum_iw_i}^{\vee}=S_{-1}^{\vee}=0$ である．
\end{proof}

\begin{lemma}[ambient への一意な延長]
\label{lemma:ambient-restriction}
制限写像
\[
 H^0(\mathcal P,\wedge^2T_{\mathcal P})
 \longrightarrow H^0(X,\wedge^2(T_{\mathcal P}|_X))
\]
は同型である．したがって $X$ 上の大域的二重ベクトルは
$\mathcal P$ 上の二重ベクトルへ一意に延長する．
この段階では Jacobi 恒等式は主張しない．
\end{lemma}

\begin{proof}
$Y=\mathcal P$ または $X$ とし，
$\mathcal V_Y=\bigoplus_i\mathcal O_Y(w_i)$，$T=T_{\mathcal P}|_Y$ とおく．
Euler 列とその外積から
\[
 0\longrightarrow\mathcal O_Y\xrightarrow{E}\mathcal V_Y
 \longrightarrow T\longrightarrow0,
 \qquad
 0\longrightarrow T\longrightarrow\wedge^2\mathcal V_Y
 \longrightarrow\wedge^2T\longrightarrow0
\]
を得る．補題~\ref{lemma:weighted-vanishing} より $H^1(Y,T)=0$ である．
また両方の重み系で $w_i<m$ および $w_i+w_j<m$（$i\ne j$）が成り立つ．
よって $\mathcal V_Y$ と $\wedge^2\mathcal V_Y$ の大域切断は，
$Y=\mathcal P,X$ で同じ多項式空間になる．$H^1(Y,\mathcal O_Y)=0$ も用いると，
両辺はともに
\[
 \frac{(\wedge^2\operatorname{Der}S)_0}
 {E\wedge(\operatorname{Der}S)_0}
\]
と同一視され，制限写像は恒等写像に対応する．
最後に $\wedge^2T_X\hookrightarrow\wedge^2(T_{\mathcal P}|_X)$ を用いる．
\end{proof}

\begin{proposition}[Casimir 代表の存在と一意性]
\label{prop:casimir-representative}
各 $\pi\in H^0(X,\wedge^2T_X)$ は，重み $0$ かつ
$\iota_{\mathrm dF}\widetilde B=0$ を満たす多項式二重ベクトル
$\widetilde B$ によって一意に表される．さらに
$\widetilde B=\iota_{\mathrm dF}C$ と書け，$C$ は次の形に一意に定まる．
\begin{align}
 d=1:\quad C&=D\wedge\partial_y\wedge\partial_z,
 &D&=\sum_{i=0}^{n-1}c_i\partial_{x_i},\quad c_i\in\mathbb C,
 \label{eq:degree-one-trivector}\\
 d=2:\quad C&=P\wedge\partial_z,
 &P&=\sum_{0\leq i<j\leq n}a_{ij}\partial_{x_i}\wedge\partial_{x_j},
 \quad a_{ij}\in\mathbb C.
 \label{eq:degree-two-trivector}
\end{align}
\end{proposition}

\begin{proof}
補題~\ref{lemma:ambient-restriction} により $\pi$ を
$B\in(\wedge^2\operatorname{Der}S)_0$ で表す．
$\pi$ が $X$ に接することは，$\iota_{\mathrm dF}B$ が $X$ 上で Euler 方向に入ることを意味する．
Euler 列の核と補題~\ref{lemma:weighted-vanishing} の $H^0$ の記述から
\[
 \iota_{\mathrm dF}B=gE+FV,
 \qquad g\in S_m,\quad V\in(\operatorname{Der}S)_0
\]
と書ける．もう一度 $\mathrm dF$ で縮約し，$E(F)=mF$ を用いると
$V(F)=-mg$ である．したがって
\begin{equation}
 \widetilde B=B-\frac1mE\wedge V
 \label{eq:casimir-correction}
\end{equation}
は同じ射影的二重ベクトルを表し，
$\iota_{\mathrm dF}\widetilde B=0$ を満たす．
$g$ の多項式への持ち上げを定数倍の $F$ だけ変えても，
$V$ の変化は Euler 方向なので \eqref{eq:casimir-correction} は変わらない．

$\operatorname{Spec}R\setminus\{0\}$ は $X$ 上の $\mathbb G_m$-torsor であり滑らかである．
また $\mathrm dF=0$ なら Euler 公式から $F=0$ である．
したがって $F_{q_0},\ldots,F_{q_{N-1}}$ の共通零点は原点のみである．
これらは Cohen--Macaulay 環 $S$ の $N$ 個の斉次パラメータなので正則列をなす．
その Koszul 複体の完全性から
\[
 \widetilde B=\iota_{\mathrm dF}C,
 \qquad C\in(\wedge^3\operatorname{Der}S)_{-m}
\]
を得る．$\partial_{q_i}\wedge\partial_{q_j}\wedge\partial_{q_k}$ の係数の次数は
$w_i+w_j+w_k-m$ である．次数 $1$ ではこれが非負となる重みの組は
$(1,2,3)$ のみ，次数 $2$ では $(1,1,2)$ のみであり，係数は定数である．
これが \eqref{eq:degree-one-trivector} と \eqref{eq:degree-two-trivector} を与える．

二つの Casimir 代表の差は $E\wedge V$ である．縮約が消えることから
$mFV=V(F)E$ を得る．$E$ の係数の最大公約数は $1$ なので $F\mid V(F)$ である．
重みから $V(F)=aF$（$a\in\mathbb C$）となり，$V=(a/m)E$，したがって
$E\wedge V=0$ である．これで $\widetilde B$ の一意性が分かる．

最後に $\iota_{\mathrm dF}C=0$ なら Koszul 完全性により
$C=\iota_{\mathrm dF}Q$，$Q\in(\wedge^4\operatorname{Der}S)_{-2m}$ と書ける．
異なる四個の重みの和は次数 $1$ で高々 $7$，次数 $2$ で高々 $5$ であり，
それぞれ $2m=12,8$ より小さい．ゆえに $Q=0$，したがって $C=0$ である．
\end{proof}

特に，選んだ重み付き座標に関して
\[
 h^0(X,\wedge^2T_X)=
 \begin{cases}
 n&d=1,\\
 \binom{n+1}{2}&d=2
 \end{cases}
\]
である．ここまでの議論は Jacobi 恒等式を一切仮定していない．

\section{Jacobi 方程式と Poisson スキーム}
\label{sec:classification}

\subsection{次数 \texorpdfstring{$1$}{1} の Jacobian 括弧}

\eqref{eq:degree-one-trivector} に対応する Casimir 代表は
\begin{equation}
 \widetilde\pi_D
 =F_zD\wedge\partial_y+D(F)\partial_y\wedge\partial_z
 +F_y\partial_z\wedge D.
 \label{eq:degree-one-bivector}
\end{equation}
したがって
\[
 \{x_i,x_j\}=0,\qquad
 \{x_i,y\}=2c_iz,\qquad
 \{y,z\}=D(F),\qquad
 \{z,x_i\}=c_iF_y.
\]
$D\ne0$ のとき $D(u)=1$ となる線型座標 $u$ と，
$D(s_i)=0$ となる独立な $n-1$ 個の線型座標 $s_i$ を取る．
$s_i$ は中心的で，残る括弧は
\[
 \{u,y\}=F_z,\qquad \{y,z\}=F_u,\qquad \{z,u\}=F_y
\]
である．

\begin{lemma}[三変数 Jacobian 括弧]
\label{lemma:jacobian-bracket}
任意の可換 $\mathbb C$-代数 $B$ と任意の $G\in B[u,v,w]$ に対し，
$B$ を中心的とし
\[
 \{u,v\}=G_w,\qquad\{v,w\}=G_u,\qquad\{w,u\}=G_v
\]
で定めた括弧は Poisson 括弧である．また $G$ は Casimir である．
\end{lemma}

\begin{proof}
Jacobiator は各引数について導分なので，$(u,v,w)$ について調べればよい．その値は
\[
 G_wG_{uv}-G_vG_{uw}-G_wG_{uv}+G_uG_{vw}
 +G_vG_{uw}-G_uG_{vw}=0
\]
である．また $\{G,u\}=G_v\{v,u\}+G_w\{w,u\}=0$ であり，
残る二つの座標も同様である．
\end{proof}

これにより，次数 $1$ ではすべての二重ベクトルが Poisson であり，
その Casimir 代表も $S$ 上で Poisson である．
$D=0$ の場合は零括弧である．

\subsection{次数 \texorpdfstring{$2$}{2} の Jacobiator}

$a_{ji}=-a_{ij}$，$a_{ii}=0$ とし，$f_i=\partial f_4/\partial x_i$，
$V_i=\sum_j a_{ij}f_j$，$V=\sum_iV_i\partial_{x_i}$ とおく．
\eqref{eq:degree-two-trivector} の Casimir 代表は
\begin{equation}
 \widetilde\pi_A=2zP+V\wedge\partial_z,
 \qquad
 \{x_i,x_j\}=2za_{ij},\quad\{x_i,z\}=V_i.
 \label{eq:degree-two-bivector}
\end{equation}
Jacobiator の規約を
\[
 J(a,b,c)=\{a,\{b,c\}\}+\{b,\{c,a\}\}+\{c,\{a,b\}\}
\]
とし，反対称な四添字の式
\begin{equation}
 q_{ijk\ell}=a_{ij}a_{k\ell}-a_{ik}a_{j\ell}+a_{i\ell}a_{jk}
 \label{eq:plucker-equation}
\end{equation}
を定める．添字が重複するとこの式は $0$ である．直接計算により
\begin{equation}
 J(x_i,x_j,z)=0,\qquad
 J_{ijk}:=J(x_i,x_j,x_k)=2\sum_{\ell=0}^n q_{ijk\ell}f_\ell.
 \label{eq:jacobi-plucker-identity}
\end{equation}
第一の等式は Hessian $(f_{ij})$ の対称性による．第二の等式は
$2(a_{jk}V_i+a_{ki}V_j+a_{ij}V_k)$ を展開すれば得られる．
したがって $J$ は $x$ 方向のみの三重ベクトルであり，その係数は $z$ に依存しない．

\begin{lemma}[分岐方程式の偏微分の独立性]
\label{lemma:gradient-independence}
$f_0,\ldots,f_n\in\mathbb C[x_0,\ldots,x_n]_3$ は線型独立である．
\end{lemma}

\begin{proof}
非零の定数ベクトル場 $D$ が $D(f_4)=0$ を満たすなら，線型座標変換後，
$f_4$ は一変数に依存しなくなる．その変数に対応する座標点では $f_4$ とすべての偏微分が消え，
分岐四次超曲面の滑らかさに矛盾する．
\end{proof}

\subsection{集合としての消滅から理想の一致へ}

\begin{definition}[固定多様体の Poisson スキーム]
\label{def:poisson-scheme}
$V_X=H^0(X,\wedge^2T_X)$，$W_X=H^0(X,\wedge^3T_X)$ とする．
二次写像
\[
 Q_X:V_X\longrightarrow W_X,\qquad
 Q_X(\pi)=\tfrac12[\pi,\pi]
\]
の零スキームを $\operatorname{Pois}_{\mathrm{sch}}(X)$ と書く．
具体的には $\operatorname{Sym}(V_X^\vee)$ において
$Q_X$ の全座標関数で生成される理想で定義する．
この定義では被約化も $\operatorname{Aut}(X)$ による商も取らない．
\end{definition}

\begin{theorem}[Jacobi ideal と Pl\"ucker ideal の一致]
\label{thm:jacobi-scheme}
固定した滑らかな $X$ と上の重み付きモデルについて次が成り立つ．
\begin{enumerate}
 \item 次数 $1$ では Jacobi ideal は零理想であり，
 $\operatorname{Pois}_{\mathrm{sch}}(X)\simeq\mathbb A^n$ である．
 \item 次数 $2$ では座標 $a_{ij}$ に関する Jacobi ideal は，
 \eqref{eq:plucker-equation} で生成される Pl\"ucker ideal と等しい．したがって
 \[
 \operatorname{Pois}_{\mathrm{sch}}(X)
 \simeq\operatorname{Spec}
 \mathbb C[a_{ij}\mid0\leq i<j\leq n]/(q_{ijk\ell})
 \]
 は $\operatorname{Gr}(2,n+1)$ の Pl\"ucker 埋め込みのアフィン錐である．
 特にこのスキームは被約かつ既約である．
\end{enumerate}
これらの同型は任意の可換 $\mathbb C$-代数上の係数族にも適用される．
多様体 $X$ 自体の変形についての主張ではない．
\end{theorem}

\begin{proof}
次数 $1$ の $Q_X$ は，前小節により $V_X(\mathbb C)$ のすべての点で消える．
係数体が無限体なので二次多項式写像として恒等的に $0$ である．

次数 $2$ で $\mathcal A=\mathbb C[a_{ij}]$ とおき，
Jacobi ideal を $I_J$，Pl\"ucker ideal を $I_P$ と書く．
式~\eqref{eq:jacobi-plucker-identity} から，普遍 Jacobiator のすべての多項式係数は
$I_P$ に属する．$X$ 上への制限と射影化は線型なので，$I_J\subset I_P$ である．

逆向きを\emph{線型な係数回収}により示す．Euler 列
$0\to\mathcal O_X\xrightarrow E\mathcal V_X\to T_{\mathcal P}|_X\to0$
から，自然な単射
\begin{equation}
 \wedge^3T_X\hookrightarrow\wedge^3(T_{\mathcal P}|_X)
 \xrightarrow{\ E\wedge\ }\wedge^4\mathcal V_X
 \label{eq:euler-jacobi-injection}
\end{equation}
を得る．第二の写像は局所的な持ち上げの選び方によらず，
Euler 方向の差を外積すると $0$ になるため well-defined である．
局所的に Euler 列を分裂すればその単射性も分かる．
普遍 Jacobiator にこの写像を適用すると，
$\partial_{x_i}\wedge\partial_{x_j}\wedge\partial_{x_k}\wedge\partial_z$
の成分は符号を除き $2zJ_{ijk}$ である．
これは $R_5\otimes\mathcal A$ の元であり，
$Q_X$ の座標関数の $\mathbb C$-線型結合を係数として持つ．

$F=z^2-f_4(x)$ が $z$ に関して monic であることから，次数付きベクトル空間として
\begin{equation}
 R_5=\mathbb C[x]_5\oplus z\mathbb C[x]_3
 \label{eq:degree-five-splitting}
\end{equation}
である．第二成分の $z$ の係数を取る線型写像を適用すると，
$2J_{ijk}=4\sum_\ell q_{ijk\ell}f_\ell$ を回収できる．
補題~\ref{lemma:gradient-independence} により，
$\mathbb C[x]_3$ 上の線型汎関数 $\lambda_\ell$ を
$\lambda_\ell(f_j)=\delta_{\ell j}$ となるよう選べる．
さらに $\lambda_\ell/4$ を適用すると，符号を除き $q_{ijk\ell}$ そのものを得る．
よって各 $q_{ijk\ell}$ は $Q_X$ の座標関数の線型結合であり，
$I_P\subset I_J$ である．ここでは零点集合や根基を取る議論を用いていない．
$n=2$ の場合は四つの異なる $x$ 添字が存在せず，両方の理想は零理想となる．

Pl\"ucker ideal の素性とアフィン錐の同定は
補題~\ref{lemma:plucker-prime} による．
最後に，各同型は $\mathbb C$ 上の座標環の同型なので，任意の可換
$\mathbb C$-代数 $B$ について $B$ 値点の同型を与える．
固定した $X$ の基底拡大では大域切断と Schouten 括弧が係数拡大と両立する．
\end{proof}

\begin{lemma}[二次 Pl\"ucker ideal の素性]
\label{lemma:plucker-prime}
$M\geq3$ とし，$\mathcal B=\mathbb C[p_{ij}\mid1\leq i<j\leq M]$ とする．
二次 Pl\"ucker 関係で生成される理想 $I$ は素理想であり，
$\operatorname{Spec}(\mathcal B/I)$ は二次元部分空間の Grassmann 多様体の
Pl\"ucker アフィン錐である．
\end{lemma}

\begin{proof}
写像
\[
 \varphi:\mathcal B\longrightarrow
 \mathbb C[u_1,\ldots,u_M,v_1,\ldots,v_M],
 \qquad p_{ij}\longmapsto u_iv_j-u_jv_i
\]
を考える．Pl\"ucker 関係は右辺で消えるので $I\subset\ker\varphi$ である．
単項式 $\prod_{s=1}^r p_{i_sj_s}$ を，
$i_1\leq\cdots\leq i_r$ かつ $j_1\leq\cdots\leq j_r$
となるよう並べられるとき標準と呼ぶ．
非標準な二列は，列を辞書式に並べた後，$i<k<\ell<j$ に対する
$p_{ij}p_{k\ell}$ と書ける．関係
\[
 p_{ij}p_{k\ell}=p_{i\ell}p_{kj}-p_{ik}p_{\ell j}
\]
でこれを書き換える．右辺第一項では上段の添字列を保ったまま下段の反転数が減り，
第二項では昇順に並べた上段の添字列が辞書式に増える．
同次数の単項式は有限個なので，この操作は停止し，すべての剰余類は標準単項式の線型結合になる．

多項式環の辞書式単項式順序を
$u_1>\cdots>u_M>v_1>\cdots>v_M$ とすると，
$\varphi(p_{ij})$ の先頭単項式は $u_iv_j$ である．
標準単項式の像の先頭単項式は
$\prod_su_{i_s}\prod_sv_{j_s}$ であり，異なる標準単項式に対して異なる．
実際，昇順の二つの添字列が一意に復元されるからである．
よってこれらの像は線型独立であり，$\ker\varphi=I$ を得る．
$\mathcal B/I$ は整域に埋め込まれるので整域である．

非零点で $p_{ij}\ne0$ とすると，残る $p_{k\ell}$ は
$p_{ij}$ および $p_{ik},p_{jk}$（$k\ne i,j$）から Pl\"ucker 関係で一意に定まる．
これは階数 $2$ の二行行列の minors の標準座標表示であり，
射影化した開集合は二次元部分空間の通常の Grassmann 座標 chart となる．
これらが非零点を覆うので，射影化は $\operatorname{Gr}(2,M)$ であり，
その同次座標環が $\mathcal B/I$ である．
\end{proof}

\begin{corollary}[大きさと特異点]
\label{cor:poisson-scheme-dimension}
次数 $2$ の Poisson スキームは次元 $2n-1$ である．
$n=2$ では $\mathbb A^3$，$n\geq3$ では非零部分が滑らかで頂点のみが特異である．
頂点での Zariski 接空間の次元は $\binom{n+1}{2}$ である．
\end{corollary}

\begin{proof}
補題~\ref{lemma:plucker-prime} の $p_{ij}\ne0$ の開集合は
一つの可逆座標と $2(n-1)$ 個の自由座標を持つので滑らかで次元 $2n-1$ である．
頂点では定義式が二次式のみなので接空間は全二重ベクトル空間である．
$n\geq3$ では
$\binom{n+1}{2}-(2n-1)=(n-1)(n-2)/2>0$ である．
\end{proof}

\subsection{持ち上げ，階数，compatibility}

次数 $2$ では $P\wedge P=0$ と $\operatorname{rank}(a_{ij})\leq2$ は同値である．
これは反対称行列を二次の交代ブロックの直和に直して確認できる．
$P\ne0$ なら線型座標を選んで $P=\partial_u\wedge\partial_v$ とでき，
残る $n-1$ 個の独立な線型座標 $s_i$ は中心的である．
括弧は
\[
 \{u,v\}=F_z,\qquad\{v,z\}=F_u,\qquad\{z,u\}=F_v
\]
となる．したがって両次数の非零構造は，中心的な $n-1$ 個の線型パラメータを持つ
三変数 Jacobian 括弧として表される．

Poisson 条件を満たすパラメータに対応する Casimir 代表は，すべて $S$ 上で Poisson なので，$\mathcal P$ 上の Poisson 構造に降下する．
$F$ が Casimir であることから $X$ も Poisson 部分多様体となる．
ambient 上の一意性は補題~\ref{lemma:ambient-restriction} による．
cone 上では $F$-Casimir という正規化を課した場合の一意性であり，
正規化なしの任意の代表についての主張ではない．

非零構造の一般階数は $2$ である．実際，三次元空間の分解可能な三重ベクトルを
一形式で縮約した二重ベクトルは階数高々 $2$ であり，射影化しても階数は増えない．
命題~\ref{prop:casimir-representative} の単射性により，非零パラメータから得た切断は非零なので，
整な $X$ の一般点での階数はちょうど $2$ である．

\begin{corollary}[二つの括弧の compatibility]
\label{cor:compatibility}
次数 $1$ の任意の二つの Poisson 構造は compatible である．
次数 $2$ の $P,Q$ に対応する二つの Poisson 構造は，
$P\wedge Q=0$ のとき，かつそのときに限り compatible である．
非零の $P,Q$ が定める二次元部分空間を $U_P,U_Q$ とすれば，
この条件は $U_P\cap U_Q\ne0$ と同値である．
\end{corollary}

\begin{proof}
次数 $1$ ではすべての線型結合が Poisson なので Schouten 括弧を偏極すればよい．
次数 $2$ では定理~\ref{thm:jacobi-scheme} により
$\pi_P+t\pi_Q$ がすべての $t$ で Poisson となることと
$(P+tQ)\wedge(P+tQ)=0$ が同値である．
$P\wedge P=Q\wedge Q=0$ なので残る条件は $P\wedge Q=0$ である．
最後の主張は四つのベクトルの線型従属性である．
\end{proof}

% !TeX root = ../../poisson-del-pezzo.tex
\section{内在的な分類と symplectic leaf}
\label{sec:geometry}
\subsection{基本線束の部分線型系}
次数1の非零定数ベクトル $D$、次数2の非零可分解二重ベクトル $P$ に対し、
それぞれ
\[
 W=\ker(D:H^0(X,H)\longrightarrow\C),\qquad
 W=\{s\in H^0(X,H)\mid\iota_{\mathrm ds}P=0\}
\]
とおく。いずれも $\dim W=n-1$ である。
逆にこの次元の部分空間 $W$ から、次数1では $D$ の直線、次数2では
2次元 annihilator の最高外積 $P$ の直線が一意に回収される。

この対応の内在性を記す。$r=n-1$、$W=\langle s_1,\ldots,s_r\rangle$ とし、
$H$ の局所自明化 $s_i=f_ie$ を用いて
\begin{equation}
 \theta_W=\left(\sum_{i=1}^{r}(-1)^{i-1}f_i\,
 \mathrm df_1\wedge\cdots\wedge\widehat{\mathrm df_i}\wedge\cdots\wedge\mathrm df_r
 \right)\otimes e^{\otimes r}
 \in H^0(X,\Omega_X^{n-2}\otimes H^{\otimes(n-1)})
 \label{eq:intrinsic-form}
\end{equation}
を定める。$f_r\neq0$ 上で括弧内は
\[
 (-1)^{r-1}f_r^r\,
 \mathrm d(f_1/f_r)\wedge\cdots\wedge\mathrm d(f_{r-1}/f_r)
\]
なので、これは局所自明化によらず貼り合う。
基底変更では行列式倍される。
$H^{\otimes(n-1)}\simeq K_X^{-1}$ と
$\Omega_X^{n-2}\otimes K_X^{-1}\simeq\bigwedge^2T_X$ により二重ベクトルの直線を得る。

これが先に分類した構造であることを局所的に確かめる。
$s_r\neq0$ 上で $t_i=s_i/s_r$ とし、活性な3変数の重みを $(a,b,c)$ とする。
重みを打ち消して得た座標を $U,V,Z$ と書けば、$X$ の局所方程式は
$G(t,U,V,Z)=0$ となる。
$G_Z\neq0$ 上で随伴公式による体積形式は、符号を除いて
\[
 \eta=\frac{\mathrm dt_1\wedge\cdots\wedge\mathrm dt_{r-1}
                    \wedge\mathrm dU\wedge\mathrm dV}{G_Z}
\]
である。Jacobian 括弧の $\{U,V\}=G_Z$ を用いると
$\iota_\pi\eta=\pm\mathrm dt_1\wedge\cdots\wedge\mathrm dt_{r-1}$ となり、
\eqref{eq:intrinsic-form} と同じ直線を与える。
両者は大域切断なので稠密開集合での一致から全体で一致する。
以上で定理\ref{thm:main}(3)の Grassmann 表示を得る。

\subsection{葉と一般階数}
$n\geq3$ で非零構造に対応する $W$ を固定する。有理写像
\begin{equation}
 \rho_W:X\dashrightarrow\mathbb P(W^*)\simeq\mathbb P^{n-2},\qquad
 p\longmapsto[s_1(p):\cdots:s_{n-1}(p)]
 \label{eq:leaf-map}
\end{equation}
を考える。比 $s_i/s_{n-1}$ は定義される開集合上の Casimir である。
cone 上の非零 Jacobian 括弧は一般点で階数2を持つ。
実際、活性方向の偏微分が全て零ならその括弧も零であり、
それが恒等的に起きないことは前節の Casimir 代表の単射性から従う。
$s_i\neq0$ 上では Euler 場 $E$ は $E(s_i)=s_i\neq0$ を満たす一方、
Hamiltonian 方向は $s_i$ を消す。従って $E$ は括弧の像に入らず、
射影化後の一般階数も2のままである。

\eqref{eq:leaf-map} のファイバーで、写像が正則であり
ファイバーが滑らかで括弧の階数が2となる部分の連結成分が symplectic leaf である。
実際、この部分では Hamiltonian 分布とファイバーの接空間が一致する。
基底集合では \eqref{eq:intrinsic-form} が零なので、階数2の葉が有理写像の未定義部分を通って延長されることはない。
一般の $W$ と一般のファイバーについて、その閉包 $Y$ は
$H$ に属する $n-2$ 個の線型切断で得られる滑らかな曲面である。
次数2では $|H|$ は基点を持たないので Bertini の定理を適用する。
次数1の唯一の基点では $x_0,\ldots,x_{n-1}$ が局所座標であるため、
独立な線型切断はそこで滑らかである。他の点では同じく Bertini を用いる。
随伴公式と交点数は
\[
 K_Y=(K_X+(n-2)H)|_Y=-H|_Y,\qquad (H|_Y)^2=H^n=d
\]
を与える。従って $Y$ は同じ次数の del Pezzo 曲面である。
特殊な $W$ で全てのファイバー閉包が滑らかになるとは主張しない。
$n=2$ では $\mathbb P(W^*)$ は点であり、次節の記述になる。

\section{曲面の退化因子}
\label{sec:surfaces}
$n=2$ では $\bigwedge^2T_X\simeq-K_X=H$ であり、全ての切断が Poisson である。
非零切断 $s$ の零因子 $D_s=(s=0)$ は反標準因子である。
$X\setminus D_s$ は連結な2次元 symplectic leaf であり、$D_s$ の各点は0次元 leaf である。
連結性は、既約な滑らかな複素多様体から複素余次元1の閉部分を除いた開部分が
連結であることによる。

\subsection{次数1}
反標準系は1個の基点を持つ pencil である。
標準方程式では基点は $x_0=x_1=0$、$z^2=y^3$ で与えられる。
$y,z$ は非零で stabilizer は自明であり、$x_0,x_1$ が局所座標となる。
従って全ての反標準メンバーは基点で滑らかである。

\begin{proposition}
次数1の非零反標準因子は被約かつ既約であり、
滑らかな楕円曲線、1個の node を持つ有理曲線、1個の cusp を持つ有理曲線の
いずれかである。
\end{proposition}
\begin{proof}
$D_s=\sum m_iC_i$ と分解すると
$1=H\cdot D_s=\sum m_i(H\cdot C_i)$ は正整数の和である。
よって成分は1個、重複度も1である。
随伴公式より $p_a(D_s)=1$。特異なら正規化の種数は0で全 $\delta$ 不変量は1となる。
滑らかな曲面の被約 Cartier 因子の特異点は平面曲線特異点であり、
$\delta=1$ の場合は node または cusp である。
具体的には、$s=0$ 上の残る重み1の座標を $u$ と取ると
$z^2=y^3+au^4y+bu^6$ となる。
$u=0$ での点は滑らかで、$u\neq0$ 上では三次式の重根を調べればよい。
$4a^3+27b^2\neq0$ なら滑らか、判別式が零で $(a,b)\neq(0,0)$ なら node、
$a=b=0$ なら cusp である。これにより特異点の分類も直接確認できる。
\end{proof}

\subsection{次数2}
反標準写像 $p:X\to\mathbb P^2$ は滑らかな四次曲線
$B=(f_4=0)$ に分岐する二重被覆である。
$D_s=p^{-1}(L)$、$L=(s=0)$ である。
\begin{center}
\begin{tabular}{lp{10cm}}
\toprule
$L\cap B$ の重複度 & $D_s$ の型\\
\midrule
$1+1+1+1$ & 滑らかな楕円曲線\\
$2+1+1$ & 既約な nodal rational curve\\
$3+1$ & 既約な cuspidal rational curve\\
$2+2$ & 2本の $(-1)$ 曲線が2点で横断的に交わる\\
$4$ & 2本の $(-1)$ 曲線が1点で接する\\
\bottomrule
\end{tabular}
\end{center}
この表を証明する。重複度 $r$ の交点近くでは解析的な単元の平方根を吸収して
$z^2=t^r$ と書ける。$r=2$ は node、$r=3$ は cusp、$r=4$ は接する2枝を与える。
$f_4|_L$ が平方になるのは全ての根の重複度が偶数のときであり、
この場合に限って逆像は2成分に分解する。
各成分は $L\simeq\mathbb P^1$ と同型で $H$ との交点数は1である。
随伴公式より自己交点は $-1$。成分間の交点数は2であり、
$2+2$ の場合は2点に、$4$ の場合は1点に集中する。
固定した曲面で表の全ての型が現れると主張するものではない。

\subsection{Jacobian 表示の符号}
中心的な重み1の座標を $s$ とし、残る変数を $(u,v,w)$ とする。
重みは次数1で $(1,2,3)$、次数2で $(1,1,2)$ であり、$m=a+b+c$ である。
座標の向きを $(s,u,v,w)$ と固定すれば
\[
 \{g,h\}=\det\frac{\partial(s,F,g,h)}{\partial(s,u,v,w)}
\]
は $\{u,v\}=F_w$、$\{v,w\}=F_u$、$\{w,u\}=F_v$ を満たす。
行を $(F,s,g,h)$ の順にすると全体の符号が逆になる。
この規約を次の量子化の関係式でも用いる。

% !TeX root = ../poisson-del-pezzo.tex
\section{滑らかな族と相対 Poisson スキーム}
\label{sec:families}

本節では多様体の定義式も動かす。まず重み付き座標を固定した族を扱い、
次節で座標変換に関する降下を行う。基底スキームは $\C$ 上の任意のスキームとし、
被約性を仮定しない。微分、接束、Schouten 括弧はすべて基底に相対的に取る。

\subsection{係数空間上の普遍族}

次数1では $x_0,\ldots,x_{n-1}$、次数2では $x_0,\ldots,x_n$ を固定する。
それぞれの斉次多項式空間を用いて
\begin{align*}
 U_1&=\{(A_4,A_6)\in\C[x]_4\oplus\C[x]_6
       \mid z^2-y^3-A_4(x)y-A_6(x)=0\text{ が滑らか}\},\\
 U_2&=\{f_4\in\C[x]_4\mid (f_4=0)\subset\Pj^n\text{ が滑らか}\}
\end{align*}
とおく。これらは係数 affine 空間の非空な開部分であり、
それぞれの標準方程式で定まる普遍族を
\[
 p_d:\mathcal X_d\longrightarrow U_d
\]
と書く。これは滑らかな射影射である。
次数1の開集合の非空性は
$F=z^2-y^3-\sum_i x_i^6$、次数2では $f_4=\sum_i x_i^4$ から分かる。
次数1の普遍超曲面も ambient stack の非自明な stabilizer を避けるので、
$\mathcal X_d$ はスキームである。

定数テンソルの空間を
\[
 C_1=\bigoplus_{i=0}^{n-1}\C\partial_{x_i},\qquad
 C_2=\bigwedge^2\left(\bigoplus_{i=0}^{n}\C\partial_{x_i}\right)
\]
とする。$D\in C_1$、$P\in C_2$ から
\begin{equation}
 D\longmapsto\iota_{\dd F}(D\wedge\partial_y\wedge\partial_z),\qquad
 P\longmapsto\iota_{\dd F}(P\wedge\partial_z)
 \label{eq:universal-casimir-map}
\end{equation}
により、普遍超曲面上の相対二重ベクトルを得る。
これらは係数空間上で正則に定義される。

\subsection{大域切断の任意基底変換}

ファイバー上の次元が一定であるという事実だけを、非被約基底上の同型の代用にはしない。
ここでは、普遍的に定義された切断が各ファイバーの基底を与えることを使う。

\begin{lemma}[一様な切断による基底変換]
\label{lemma:universal-sections-base-change}
$p:Y\to S$ を平坦かつ有限表示の固有射、$\mathcal E$ を $Y$ 上の有限局所自由層、
$\mathcal V$ を $S$ 上の有限局所自由層とする。
写像 $p^*\mathcal V\to\mathcal E$ がすべての幾何学的点 $s$ で同型
\[
 \mathcal V_s\xrightarrow{\sim}H^0(Y_s,\mathcal E_s)
\]
を誘導すると仮定する。このとき随伴写像
$\mathcal V\to p_*\mathcal E$ は同型であり、任意の $g:T\to S$ に対し
\[
 g^*\mathcal V\xrightarrow{\sim}(p_T)_*\mathcal E_T
\]
である。特に $p_*\mathcal E$ の形成は任意の基底変換と可換である。
\end{lemma}

\begin{proof}
固有平坦射に対する perfect base change
\cite[Tag 0B91]{StacksProject} により、$Rp_*\mathcal E$ は perfect であり、
その導来引き戻しは任意の基底変換後の $R(p_T)_*\mathcal E_T$ に等しい。
これらの複体には負次数のコホモロジーがないので、$S$ 上局所的に
$Rp_*\mathcal E$ を非負次数の有限自由複体
\[
 K^0\xrightarrow{d^0}K^1\longrightarrow\cdots\longrightarrow K^b
\]
で表せる。実際、任意の bounded finite free 代表の負の最低次数から始め、
ファイバーで完全な直和因子を小行列式の可逆化によって順に消去すればよい。

$\mathcal V$ を自明化し、その基底の像を $K^0$ の cycle
$s_1,\ldots,s_r$ で表す。固定した点 $s$ でこれらは
$\ker(d^0\otimes\kappa(s))$ の基底である。
従って近傍を縮めれば、$s_i$ を $K^0$ の基底の一部に延長して
\[
 K^0=\langle s_1,\ldots,s_r\rangle\oplus L
\]
と書ける。$d^0|_L$ は $s$ の剰余体上で単射なので、さらに小行列式を可逆化すれば
分裂単射になる。これは任意の基底変換後も単射である。
一方 $d^0(s_i)=0$ なので、任意の $T\to S$ について
\[
 \ker(d^0\otimes\mathcal O_T)
 =\langle s_1,\ldots,s_r\rangle\otimes\mathcal O_T
\]
を得る。導来基底変換の次数0を取れば結論となる。
この議論では $S$ や $T$ の被約性を用いていない。
\end{proof}

\begin{proposition}[相対二重ベクトル束]
\label{prop:relative-bivectors}
式~\eqref{eq:universal-casimir-map} は同型
\begin{equation}
 \mathcal O_{U_d}\otimes C_d
 \xrightarrow{\sim}(p_d)_*\bigwedge^2T_{\mathcal X_d/U_d}
 \label{eq:relative-bivector-bundle}
\end{equation}
を与え、その形成は任意の基底変換と可換である。
従って任意の $T\to U_d$ 上の相対大域二重ベクトルは、
$C_d\otimes\mathcal O_T$ の切断に一意に対応する。
\end{proposition}

\begin{proof}
命題~\ref{prop:casimir-representative} の証明は任意の代数閉な標数0の拡大体上で同じなので、
式~\eqref{eq:universal-casimir-map} は各幾何学的ファイバーで大域切断の同型を与える。
$\mathcal X_d\to U_d$ は滑らかな射影射であり、$\bigwedge^2T_{\mathcal X_d/U_d}$ は
有限局所自由である。したがって補題~\ref{lemma:universal-sections-base-change} を適用できる。
\end{proof}

ここで必要なのは $\bigwedge^2T$ の大域切断の同定であり、
$(p_d)_*\bigwedge^3T_{\mathcal X_d/U_d}$ が局所自由であるとの仮定は置かない。
次の構成では相対 Euler 列を用いて Jacobiator の消滅を有限個の係数で検出する。

\subsection{相対 Jacobi イデアル}

各 $T\to U_d$ に対し
\[
 \underline{\Pois}_{\mathcal X_d/U_d}(T)=
 \left\{\pi\in H^0(\mathcal X_d\times_{U_d}T,
                 \bigwedge^2T_{\mathcal X_d\times_{U_d}T/T})
                \ \middle|\ [\pi,\pi]=0\right\}
\]
と定める。この等式は $T$ 上の切断としての等式であり、
幾何学的ファイバー上でのみ Jacobi 恒等式を要求するものではない。
$\mathfrak C_{2,n}$ を二次 Pl\"ucker 関係で定義される affine cone
\[
 \mathfrak C_{2,n}=
 \operatorname{Spec}\C[a_{ij}\mid0\leq i<j\leq n]/(q_{ijk\ell})
\]
と書き、$\mathfrak C_{1,n}=\A^n$ とおく。
ここで錐は Pl\"ucker 方程式で定義したスキームを意味する。
任意の係数環上の点を必ず二つのベクトルの外積として書ける，という意味ではない。
例えば $n\geq3$ の双対数上では，
$\varepsilon(e_0\wedge e_1+e_2\wedge e_3)$ もこの錐の点となる。
このような無限小点を含めることが，以下の基底変換と二次障害の記述に必要である。

\begin{theorem}[相対 Poisson スキーム]
\label{thm:relative-poisson}
上の関手は $U_d$ 上の閉 affine スキームで表現され、
標準座標の下で同型
\begin{equation}
 \mathcal Z_{d,n}:=\underline{\Pois}_{\mathcal X_d/U_d}
 \simeq U_d\times\mathfrak C_{d,n}
 \label{eq:relative-poisson-product}
\end{equation}
がある。この同型は非被約基底を含む任意の基底変換と両立する。
特に次数2の相対 Jacobi イデアルそのものが、
定義式の係数に依存しない Pl\"ucker イデアルに等しい。
\end{theorem}

\begin{proof}
命題~\ref{prop:relative-bivectors} により相対二重ベクトル全体は
$U_d\times C_d$ で表される。
次数1では $D,\partial_y,\partial_z$ は互いに可換な導分であり、
式~\eqref{eq:degree-one-bivector} はこれらに関する Jacobian 括弧である。
補題~\ref{lemma:jacobian-bracket} と同じ混合偏微分の相殺は、
三つの可換導分に対してもそのまま成り立つ。
従って任意の係数環で Jacobiator が零となる。
ここでは $D$ のいずれの係数も可逆化していない。

次数2では $U_2$ の affine 開部分 $\operatorname{Spec}A$ で考える。
普遍方程式を $F=z^2-f_4(x)$ とし、$B=A[a_{ij}]$ とおく。
相対 Euler 列と法線束列は局所分裂するベクトル束の完全列なので、
式~\eqref{eq:euler-jacobi-injection} の相対版
\begin{equation}
 \bigwedge^3T_{\mathcal X_2/U_2}
 \hookrightarrow\bigwedge^4
       \left(\bigoplus_i\mathcal O_{\mathcal X_2}(w_i)\right)
 \label{eq:relative-jacobi-injection}
\end{equation}
は任意の基底変換後も単射である。
また monic な式 $F=z^2-f_4$ により、任意の係数 $A$-代数 $B'$ に対し
\[
 \bigl(B'[x,z]/(F)\bigr)_t
   =B'[x]_t\oplus zB'[x]_{t-2}
\]
が成り立つ。これらは有限自由 $B'$-加群である。
補題~\ref{lemma:weighted-vanishing} の \v Cech 計算は任意の係数環で有効なので、
この次数付き加群は $\mathcal O_{\mathcal X_{2,B'}}(t)$ の大域切断と一致する。
従って普遍 Jacobiator の
\eqref{eq:relative-jacobi-injection} による像の全単項式係数は
$B$ の有限個の二次多項式を与え、その零スキームが求める関手を表す。
それらの生成するイデアルを $I_J$ とする。

式~\eqref{eq:jacobi-plucker-identity} は $B[x,z]$ でも成立するので、
Pl\"ucker イデアルを $I_P$ とすると $I_J\subset I_P$ である。
逆向きは係数回収を相対化して示す。
普遍偏微分 $f_j=\partial f_4/\partial x_j$ の三次単項式係数を列に持つ行列を
$M(f)$ とする。補題~\ref{lemma:gradient-independence} により
すべての幾何学的点で $M(f)$ の階数は $n+1$ である。
従って $(n+1)\times(n+1)$ 小行列式 $\Delta$ が可逆な開集合が $U_2$ を覆う。
その上では可逆小行列の逆行列から $A$-線型写像
\[
 \lambda_\ell:A[x]_3\longrightarrow A,
 \qquad\lambda_\ell(f_j)=\delta_{\ell j}
\]
を作れる。この写像は任意の $A$-代数への基底変換後も同じ性質を持つ。

$E\wedge J$ の
$\partial_{x_i}\wedge\partial_{x_j}\wedge\partial_{x_k}\wedge\partial_z$
成分は、成分を並べる符号を除き $2zJ_{ijk}$ である。
その $zA[x]_3$ 成分から $z$ を除き、$\lambda_\ell/4$ を適用すると、
式~\eqref{eq:jacobi-plucker-identity} によって $q_{ijk\ell}$ が得られる。
したがって各 $q_{ijk\ell}$ は Jacobiator の係数の $A$-線型結合であり、
$I_P\subset I_J$ となる。
この回収は $\Delta$ を可逆化した開集合ごとに成立し、イデアルの等式として貼り合う。
任意の基底変換にも同じ左逆を引き戻せるので、非被約基底でも等式は保たれる。
$n=2$ では四つの異なる $x$ 添字がなく、両イデアルは零である。
\end{proof}

\begin{corollary}[平坦性と相対的な特異点]
\label{cor:relative-flatness}
$\mathcal Z_{d,n}\to U_d$ は affine、有限表示、平坦であり、
幾何学的ファイバーは被約かつ既約である。
次数1では相対次元 $n$ の滑らかな射である。
次数2では相対次元 $2n-1$ であり、$n=2$ なら滑らかである。
$n\geq3$ の次数2では非零部分上で滑らかで、
非滑らか点の集合は零切断に一致する。
\end{corollary}

\begin{proof}
積表示~\eqref{eq:relative-poisson-product} と
補題~\ref{lemma:plucker-prime}、系~\ref{cor:poisson-scheme-dimension} を用いる。
平坦性は $\mathcal Z_{d,n}\to U_d$ が $\mathfrak C_{d,n}\to\operatorname{Spec}\C$
の基底変換であることによる。
非滑らか点についての主張は点集合について述べており、
Jacobian イデアルで定めた特異部分スキームが被約であるとの主張ではない。
\end{proof}

\subsection{多様体を動かしたときの二次障害}

\begin{corollary}
\label{cor:relative-deformations}
次数2で $u\in U_2(\C)$ を固定する。
零 Poisson 構造における座標付き同時変形の一次空間は
\[
 T_uU_2\oplus\bigwedge^2\C^{n+1}
\]
である。$\varepsilon^2=0$ 上の一次変形を $(\dot f,Q)$ と書くと、
これを $\varepsilon^3=0$ 上へ持ち上げられる必要十分条件は
\[
 Q\wedge Q=0
\]
である。この条件は定義式の一次変形 $\dot f$ に依存せず、
多様体の変形を加えても非零の二次障害を相殺できない。
条件が成立すれば、この一次変形は $\C[[t]]$ 上へ持ち上がる。
\end{corollary}

\begin{proof}
積表示により、二重ベクトルパラメータを
$P=\varepsilon Q+\varepsilon^2R$ と書くと
\[
 P\wedge P=\varepsilon^2Q\wedge Q\pmod{\varepsilon^3}.
\]
定義式の係数はこの式に現れない。
$U_2$ は affine 空間の開部分なので、定義式の任意の一次方向を
$f(t)=f_4+t\dot f$ で形式的に延長できる。
$Q\wedge Q=0$ の場合、$P(t)=tQ$ と組み合わせれば全次数への延長を得る。
\end{proof}

この系は座標付きの同時変形についての主張である。
座標変換による商の接空間、安定化群、極性の線束のスカラー自己同型は次節で扱う。
また、相対 Poisson スキームの積表示は選んだ標準座標上の表示であり、
座標を忘れたモジュライ上で自明な積になることを意味しない。

% !TeX root = ../../poisson-del-pezzo.tex
\section{座標変換と Poisson モジュライ}
\label{sec:moduli}

前節の積表示は，方程式と重み付き座標を選んだ上での表示である．
この節では座標変換の作用を計算し，次数 $1,2$ の滑らかな del Pezzo manifold と
Poisson 構造の族を分類する代数的スタックを構成する．
偏極線束のスカラー自己同型と Poisson 構造のスカラー倍は，異なる操作として扱う．

\subsection{偏極を持つ族と相対切断環}

$T$ を $\C$ 上のスキームとする．\emph{線束を指定した del Pezzo 族}とは，
滑らかで射影的な相対次元 $n$ の射 $p:\mathcal X\to T$ と相対 ample 線束
$\mathcal H$ の組で，各幾何学的ファイバーが連結であり，
\[
 \omega_{X_t}^{-1}\simeq H_t^{\otimes(n-1)},\qquad H_t^n=d
\]
を満たすものをいう．ここでは反標準線束との同型そのものは指定しない．
射は族の同型と偏極線束の同型の組であり，線束のスカラー自己同型を含む．
このスタックを $\widetilde{\mathcal D}_{n,d}$ と書く．
相対 Poisson 二重ベクトル $\pi\in\Gamma(\mathcal X,\wedge^2T_{\mathcal X/T})$
をさらに指定したスタックを $\widetilde{\mathcal P}_{n,d}$ と書く．

一方，\emph{偏極類を指定した族}のスタック
$\mathcal D_{n,d}$ は，このスカラー $\mathbb G_m$ を自己同型群から除いた
fppf スタックとする．より具体的には，線束を局所的に選び，
基底からの線束による $\mathcal H\mapsto\mathcal H\otimes p^*M$ を許して
降下させた偏極類を指定する．対応する Poisson スタックを $\mathcal P_{n,d}$ と書く．
従って $\mathcal D_{n,d}$ 上に基本線束そのものが大域的に存在することは，定義に含めない．
この区別は，任意の基底上の族と自己同型群を正しく記述するために必要である．

\begin{lemma}[基本偏極類の一意性と無標識の族]
\label{lemma:fundamental-polarization-class}
$p:\mathcal X\to T$ を滑らかで射影的な族とし，各幾何学的ファイバーが
$n$ 次元，次数 $d$ の del Pezzo manifold であるとする．
このとき相対 Picard 層 $\operatorname{Pic}_{\mathcal X/T}$ に
\[
 (n-1)h=[\omega_{\mathcal X/T}^{-1}]
\]
を満たす切断 $h$ がただ一つ存在する．これは fppf 局所的に基本線束で表される．
従って $\mathcal D_{n,d}$ は，基本偏極を別途指定しない滑らかな del Pezzo 族の
スタックとも同値である．Poisson 構造を付けても同じ結論が成り立つ．
\end{lemma}

\begin{proof}
Fano ファイバーの Kodaira 消滅から，任意基底変換と両立する
$Rp_*\mathcal O_{\mathcal X}\simeq\mathcal O_T$ がある．
相対 Picard 層は $T$ 上局所有限表示の代数空間である
\cite[Tags 0D2C, 0DNI]{StacksProject}．ここではさらに，それが $T$ 上 \'{e}tale であることを示す．
アフィン基底の平方零拡大 $T_0\subset T_1$ とそのイデアル $J$ に対し，
$p_{T_1}$ は滑らかな全射なので，$T_1$ の \'{e}tale 被覆上で切断 $\sigma$ を持つ．
必要ならその被覆を affine 開集合に細分する．
切断を持つ Picard 関手では，$T_0$ 上の Picard 類を
$\sigma$ に沿って自明化した実際の線束で表せる
\cite[Tag 0D28]{StacksProject}．そこで
単元の完全列
\[
 1\longrightarrow1+p^*J\longrightarrow
 \mathcal O_{\mathcal X_{T_1}}^\times\longrightarrow
 \mathcal O_{\mathcal X_{T_0}}^\times\longrightarrow1
\]
を用いる．$1+p^*J\simeq p^*J$ であり，導来射影公式から
$H^1(\mathcal X_{T_0},p^*J)=H^2(\mathcal X_{T_0},p^*J)=0$ である．
従って線束の持ち上げは存在し，その同型類は一意である．
この一意性によって上の \'{e}tale 被覆上の Picard 類の持ち上げが降下し，
相対 Picard 層の形式的 \'{e}tale 性を得る．
局所有限表示性と合わせて \'{e}tale 性を得る．

幾何学的ファイバーの Picard 群は torsion-free である．
実際，torsion 線束 $M$ に対し $M\otimes\omega_X^{-1}$ は ample なので
Kodaira 消滅が $H^i(X,M)=0$ ($i>0$) を与え，
Riemann--Roch によって $\chi(M)=\chi(\mathcal O_X)=1$ となる．
従って $M$ は非零切断を持つが，その零因子は数値的に自明であるため
ample 因子との交点数から空であり，$M\simeq\mathcal O_X$ である．

$r=n-1$ とし，$[r]:\operatorname{Pic}_{\mathcal X/T}\to
\operatorname{Pic}_{\mathcal X/T}$ を考える．
これは $T$ 上 \'{e}tale な代数空間の間の $T$-射なので \'{e}tale である．
$[\omega^{-1}]$ に沿って引き戻した根の空間は，各幾何学的ファイバーに
仮定から少なくとも一点を持ち，torsion-free 性からちょうど一点を持つ．
従ってこの根の空間は $T$ への全射かつ radicial な \'{e}tale 射であり，同型である．
これが求める切断 $h$ である．
その fppf 局所的な線束表示は Picard 層の定義
\cite[Tag 0DME]{StacksProject} により存在し，ファイバー上で ample であるから相対 ample である．
一意性により，族のすべての同型はこの偏極類を保つ．
線束そのものを大域的に選べるとは結論していない．
\end{proof}

\begin{lemma}[族の標準形]
\label{lemma:family-normal-form}
線束を指定した del Pezzo 族は，基底の fppf 被覆上で
\eqref{eq:degree-one-model} または \eqref{eq:degree-two-model} の族に同型である．
その際，$\mathcal H$ は $\mathcal O_{\mathcal X}(1)$ に対応する．
この主張は非被約基底にも成り立つ．
\end{lemma}

\begin{proof}
各ファイバーの標準モデルと Kodaira 消滅から，$r\geq0$ に対し
$H^i(X_t,H_t^r)=0$ ($i>0$) であり，$h^0(X_t,H_t^r)$ は $n,d,r$ のみで決まる．
従って相対切断環
\[
 \mathcal R=\bigoplus_{r\geq0}p_*\mathcal H^r
\]
の各次数部分は局所自由で，任意の基底変換と両立する．
ここでは高次コホモロジーが消える形の cohomology and base change を使っており，
基底の被約性は仮定していない．

次数 $2$ では $\mathcal R_1$ の基底を $x_0,\ldots,x_n$ とし，
$\mathcal R_2/\operatorname{Sym}^2\mathcal R_1$ の局所的な生成元を $z$ に持ち上げる．
この商は線束である．実際，$\operatorname{Sym}^2\mathcal R_1\to\mathcal R_2$
は各ファイバーで単射なので，局所的に分裂単射になる．
次数 $1$ では $\mathcal R_1$ の基底 $x_0,\ldots,x_{n-1}$ を選び，
\[
 \mathcal R_2/\operatorname{Sym}^2\mathcal R_1,
 \qquad
 \mathcal R_3/(\mathcal R_1\mathcal R_2)
\]
という二つの線束の生成元をそれぞれ $y,z$ に持ち上げる．
第二の商については，$y$ を選んだ後の
$\operatorname{Sym}^3\mathcal R_1\oplus\mathcal R_1y\to\mathcal R_3$ が
各ファイバーで単射なので分裂単射であることを使う．
これらは重み付き多項式環から $\mathcal R$ への次数付き全射を与える．
実際，各次数の余核は有限型で，すべての幾何学的ファイバー上で零なので零である．

最初の関係は次数 $m=4$ または $6$ の線束をなし，
局所的な生成元 $F$ の $z^2$ 係数は可逆である．
その係数を $1$ とし，$z$ の平方完成を行う．
次数 $1$ では $y^3$ 係数も可逆であるため，fppf 局所的に単元の三乗根を取り，
さらに $y$ の平行移動で $y^2$ 項を消すと所定の標準形を得る．
次数 $2$ では平方完成だけで標準形になる．
標準形の $F$ は $z$ に関して monic なので，$S/(F)$ の各次数部分は局所自由であり，
その階数は $\mathcal R$ の対応する部分と一致する．
上の全射 $S/(F)\to\mathcal R$ は各次数で同階数の局所自由加群の全射なので同型である．
相対 ample 線束の切断環から $\mathcal X$ を $\operatorname{Proj}_T\mathcal R$ として回収すればよい．
\end{proof}

\subsection{正規形を保つ群と二重ベクトルの作用}

次数 $1$ では
\[
 \widetilde G_1=\operatorname{GL}_n\times\mathbb G_m,
 \qquad
 (A,t):(x,y,z)\longmapsto(Ax,t^2y,t^3z)
\]
を取り，次数 $2$ では
\[
 \widetilde G_2=\operatorname{GL}_{n+1}\times\mathbb G_m,
 \qquad
 (A,b):(x,z)\longmapsto(Ax,bz)
\]
を取る．これらは以下の作用で滑らかな係数空間 $U_d$ を保つ．

\begin{proposition}[座標・方程式・Poisson パラメータの同時変換]
\label{prop:coordinate-action}
標準形間のすべての次数付き切断環の同型は，基底上局所的に上記の形になる．
係数と Casimir パラメータへの作用は
\begin{align}
 d=1:\quad
 (A_4,A_6,D)&\longmapsto
 \bigl(t^4 A_4(A^{-1}x),\ t^6 A_6(A^{-1}x),\ t^{-1}AD\bigr),
 \label{eq:moduli-action-one}\\
 d=2:\quad
 (f_4,P)&\longmapsto
 \bigl(b^2f_4(A^{-1}x),\ b^{-1}(\wedge^2A)P\bigr)
 \label{eq:moduli-action-two}
\end{align}
である．重み付きスカラー
\[
 Z_1=\{(\lambda I_n,\lambda)\},\qquad
 Z_2=\{(\lambda I_{n+1},\lambda^2)\}
\]
は係数と Poisson パラメータの両方に自明に作用する．
従って作用は滑らかな代数群
\begin{equation}
 G_1=\widetilde G_1/Z_1\simeq\operatorname{GL}_n,
 \qquad
 G_2=\widetilde G_2/Z_2
 \label{eq:moduli-groups}
\end{equation}
に降下する．
\end{proposition}

\begin{proof}
次数 $2$ の次数付き同型は，関係式より低い次数の生成元を比較すると
$x'=Ax$，$z'=bz+q_2(x)$ と書ける．
関係式の $z$ に関する一次係数を比較して $2bq_2=0$，従って $q_2=0$ を得る．
次数 $1$ では一般に
\[
 x'=Ax,\qquad y'=ay+q_2(x),\qquad
 z'=bz+y\ell_1(x)+q_3(x)
\]
である．関係式の $z$ の一次係数から $\ell_1=q_3=0$，
続いて $y^2$ 係数から $q_2=0$ を得る．
$z^2,y^3$ 係数の比較は $b^2=a^3$ を与え，$t=b/a$ とすれば $a=t^2,b=t^3$ となる．
ここでは $2,3$ および $a,b$ が可逆であることだけを用いるので，非被約基底でも同じ議論が成り立つ．

座標変換を $\varphi$，変換後の標準方程式を $F'$ と書く．
$F'\circ\varphi=cF$ なら，縮約の自然性から
\begin{equation}
 \varphi_*\bigl(\iota_{\dd F}C\bigr)
 =\iota_{\dd F'}\bigl(c^{-1}\varphi_*C\bigr)
 \label{eq:equation-line-covariance}
\end{equation}
となる．次数 $1$ では $c=t^6$，
$\varphi_*(D\wedge\partial_y\wedge\partial_z)
=t^5(AD)\wedge\partial_{y'}\wedge\partial_{z'}$ である．
次数 $2$ では $c=b^2$，
$\varphi_*(P\wedge\partial_z)=b(\wedge^2A)P\wedge\partial_{z'}$ である．
これで \eqref{eq:moduli-action-one}，\eqref{eq:moduli-action-two} を得る．
中央部分群の作用は代入で確認できる．
$G_1$ の同型は $(A,t)\mapsto t^{-1}A$ により与えられる．
\end{proof}

\begin{remark}[二重被覆の involution]
\label{rem:deck-action}
$\tau:z\mapsto-z$ は次数 $1$ では $(A,t)=(I,-1)$，
次数 $2$ では $(A,b)=(I,-1)$ である．いずれも方程式と $x$ の一次切断を保つが，
Poisson パラメータには $-1$ で作用する．
従って二重ベクトルの affine パラメータを，基本線束の一次切断の作用だけから
同変に復元することはできない．式~\eqref{eq:equation-line-covariance} の
方程式の変換因子が必要である．
\end{remark}

\subsection{内在的な線束の捻り}

前項の変換因子を，座標を使わずに記述する．
線束を指定した族に対し
\begin{equation}
 \mathcal E=p_*\mathcal H,\qquad
 \mathcal L=p_*\bigl(\omega_{\mathcal X/T}^{-1}
                   \otimes\mathcal H^{-(n-1)}\bigr)
 \label{eq:moduli-intrinsic-bundles}
\end{equation}
とおく．$\mathcal E$ の階数は次数 $1$ で $n$，次数 $2$ で $n+1$ である．
$\mathcal L$ は線束であり，評価写像は同型
\[
 p^*\mathcal L\simeq
 \omega_{\mathcal X/T}^{-1}\otimes\mathcal H^{-(n-1)}
\]
を与える．実際，右辺は各ファイバーで自明，$h^0=1$，$H^1=0$ なので
cohomology and base change と評価写像のファイバー上の同型性を適用できる．

\begin{proposition}[内在的な二重ベクトル束と Poisson 錐]
\label{prop:intrinsic-bivector-bundle}
任意の線束を指定した del Pezzo 族について，自然な同型
\begin{equation}
 p_*\wedge^2T_{\mathcal X/T}
 \simeq\bigwedge^{n-1}\mathcal E\otimes\mathcal L
 \label{eq:intrinsic-bivector-bundle}
\end{equation}
があり，任意の基底変換と両立する．
次数 $1$ の相対 Poisson スキームはこのベクトル束の全空間である．
次数 $2$ では，この束の中の分解可能な外積の Pl\"ucker 錐である．
いずれも $\mathcal H$ を基底の線束で捻る操作に不変であり，
偏極類のスタック $\mathcal D_{n,d}$ 上へ降下する．
\end{proposition}

\begin{proof}
式~\eqref{eq:intrinsic-form} を相対微分で書けば，交代多重線型な写像
\[
 \bigwedge^{n-1}\mathcal E
 \longrightarrow
 p_*\bigl(\Omega_{\mathcal X/T}^{n-2}\otimes\mathcal H^{n-1}\bigr)
\]
を得る．これに $\mathcal L$ をテンソルし，評価同型と
$\Omega_{\mathcal X/T}^{n-2}\otimes\omega_{\mathcal X/T}^{-1}
\simeq\wedge^2T_{\mathcal X/T}$ を用いる．
ファイバー上では，標準座標の外積基底は，補集合の添字を持つ
$D$ または $P$ の基底に共通の非零因子と符号を除いて対応する．
従ってこの写像はファイバーごとに同型である．
命題~\ref{prop:relative-bivectors} と標準形の fppf 局所存在により右辺も局所自由で
任意の基底変換と両立するため，写像は任意の基底上で同型である．

標準座標上ではこの外積同型は，次数 $2$ の二次 Pl\"ucker 関係を
$(n-1)$ 次外積の Pl\"ucker 関係に移す．
定理~\ref{thm:relative-poisson} により，これは閉点だけでなく閉部分スキームの同定である．
次数 $1$ では階数 $n$ の束の $(n-1)$ 次外積なので Pl\"ucker 関係はない．
最後に，$\mathcal H\mapsto\mathcal H\otimes p^*M$ に対して
\[
 \mathcal E\mapsto\mathcal E\otimes M,
 \qquad \mathcal L\mapsto\mathcal L\otimes M^{-(n-1)}
\]
となり，\eqref{eq:intrinsic-bivector-bundle} の二つの因子は相殺する．
この同型は評価と微分の自然な写像から作られているので，降下の cocycle 条件も満たす．
\end{proof}

\subsection{商スタックの構成と幾何学}

\[
 \mathfrak C_{1,n}=\mathbb A^n,\qquad
 \mathfrak C_{2,n}=\operatorname{Spec}
 \C[a_{ij}\mid0\leq i<j\leq n]/(q_{ijk\ell})
\]
とおく．以下で用いる $G_d$ の作用は
\eqref{eq:moduli-action-one}，\eqref{eq:moduli-action-two} の作用である．

\begin{lemma}[無限小自己同型の消滅]
\label{lemma:no-vector-fields}
次数 $1,2$ の滑らかな del Pezzo manifold は $H^0(X,T_X)=0$ を満たす．
従って $G_d$ における標準方程式の幾何学的 stabilizer は有限で被約である．
\end{lemma}

\begin{proof}
Euler 列と $H^1(X,\mathcal O_X)=0$ により，大域ベクトル場は
$V\in(\operatorname{Der}S)_0$ で表される．
$w_i<m$ なので各係数は $S$ に一意に持ち上がり，
接条件から $V(F)=cF$ ($c\in\C$) となる．
$V$ を $V-(c/m)E$ に替えると $V(F)=0$ とできる．
偏微分列の Koszul 完全性により，この syzygy は係数の次数が
$w_i+w_j-m$ の Koszul 関係の線型結合である．
異なる $i,j$ について常に $w_i+w_j<m$ なので，すべての係数は零となる．
従ってもとの $V$ は Euler 方向だけであり，$X$ 上では零である．
命題~\ref{prop:coordinate-action} により，$G_d$ の stabilizer は偏極類を保つ
自己同型群である．その Lie 代数は零で，有限型の代数群だから有限かつ被約である．
\end{proof}

\begin{theorem}[Poisson 構造付き del Pezzo manifold のモジュライ]
\label{thm:poisson-moduli}
$n\geq2$，$d\in\{1,2\}$ とする．次の自然な同型がある．
\begin{align}
 \widetilde{\mathcal D}_{n,d}&\simeq[U_d/\widetilde G_d],&
 \widetilde{\mathcal P}_{n,d}&\simeq[(U_d\times \mathfrak C_{d,n})/\widetilde G_d],
 \label{eq:unrigidified-moduli}\\
 \mathcal D_{n,d}&\simeq[U_d/G_d],&
 \mathcal P_{n,d}&\simeq[(U_d\times \mathfrak C_{d,n})/G_d].
 \label{eq:poisson-moduli-quotient}
\end{align}
$\mathcal D_{n,d}$ と $\mathcal P_{n,d}$ は有限型の既約な
Deligne--Mumford スタックである．$\mathcal D_{n,d}$ は滑らかであり，
忘却射
\begin{equation}
 \mathcal P_{n,d}\longrightarrow\mathcal D_{n,d}
 \label{eq:poisson-forgetful}
\end{equation}
は affine，有限表示かつ平坦で，零切断を持つ．
その幾何学的ファイバーは次数 $1$ で $\mathbb A^n$，次数 $2$ で
$\operatorname{Gr}(2,n+1)$ の Pl\"ucker 錐である．
すべての構成は非被約基底を含む任意の基底変換と両立する．
\end{theorem}

\begin{proof}
補題~\ref{lemma:family-normal-form} が局所的な対象の全射性を，
命題~\ref{prop:coordinate-action} が任意の基底上の同型の記述を与える．
従って標準形の同型の fppf 層は $\widetilde G_d$-torsor となり，
最初の商表示を得る．Poisson 構造については
定理~\ref{thm:relative-poisson} の相対分類をこの torsor 上で用い，
命題~\ref{prop:coordinate-action} に従って降下させる．
中央部分群 $Z_d$ は切断環の重み付きスカラーであり，
射影化した族上で恒等，偏極線束上でスカラー自己同型となる．
これを除くことは $\widetilde G_d$ を $G_d$ に置き換えることに一致する．

$U_d$ と $\mathfrak C_{d,n}$ は有限型で，$G_d$ は滑らかな有限型代数群だから，
商は有限型代数的スタックである \cite[Tag 06FI]{StacksProject}．
補題~\ref{lemma:no-vector-fields} によりすべての幾何学的 stabilizer は unramified である．
Poisson 構造付きの stabilizer はその閉部分群なので，同じ結論が成り立つ．
従って両方の商は Deligne--Mumford である
\cite[Tags 06N3, 0DSL]{StacksProject}．
$U_d$ は非空な affine 空間の開集合なので滑らかかつ既約であり，
$\mathfrak C_{d,n}$ は被約かつ既約である．商の既約性と基底スタックの滑らかさはその atlas から従う．
非空性は，例えば次数 $1$ では $F=z^2-y^3-\sum x_i^6$，
次数 $2$ では $F=z^2-\sum x_i^4$ から分かる．

忘却射を atlas $U_d\to\mathcal D_{n,d}$ で引き戻すと
$U_d\times \mathfrak C_{d,n}\to U_d$ になる．
affine 性，有限表示性，平坦性および記載したファイバーは，この表示と fppf 降下から従う．
零点は群作用で保たれるので零切断も降下する．
任意基底変換との両立は相対分類と torsor の降下から従う．
\end{proof}

\begin{remark}[固定対象の fiber と自己同型商]
族 $(X,H)$ を表す $\operatorname{Spec}\C\to\mathcal D_{n,d}$ に沿う
\eqref{eq:poisson-forgetful} のファイバーは，$X$ との同型を記憶する
$\Pois_{\mathrm{sch}}(X)$ そのものである．
一方，$X$ の同型類を固定しつつその自己同型も許すスタックは
$[\Pois_{\mathrm{sch}}(X)/\operatorname{Aut}(X,[H])]$ である．
例えば注意~\ref{rem:deck-action} の involution は $\pi$ と $-\pi$ を同型にするが，
非零 $\pi$ 自身の stabilizer には入らない．
\end{remark}

\begin{corollary}[次元と滑らかな部分]
\label{cor:moduli-dimension}
\begin{align*}
 b_{n,1}&=\binom{n+3}{4}+\binom{n+5}{6}-n^2,&
 \dim\mathcal P_{n,1}&=b_{n,1}+n,\\
 b_{n,2}&=\binom{n+4}{4}-(n+1)^2,&
 \dim\mathcal P_{n,2}&=b_{n,2}+2n-1
\end{align*}
とおくと，$\dim\mathcal D_{n,d}=b_{n,d}$ である．
$\mathcal P_{n,1}$ は滑らかである．
$\mathcal P_{2,2}$ も滑らかであり，$n\geq3$ の $\mathcal P_{n,2}$ の
非滑らか点の集合は零切断に一致する．
特に，曲面の場合の $\mathcal P_{2,1},\mathcal P_{2,2}$ の次元は $10,9$，
三次元の場合の $\mathcal P_{3,1},\mathcal P_{3,2}$ の次元は $37,24$ である．
\end{corollary}

\begin{proof}
係数空間の次元から $\dim G_1=n^2$，$\dim G_2=(n+1)^2$ を引き，
Poisson ファイバーの次元を加える．
滑らかさについては atlas $U_d\times \mathfrak C_{d,n}$ と
系~\ref{cor:poisson-scheme-dimension} を用いる．
ここで主張するのはスタックの非滑らか点の集合であり，
特異点スキームの scheme structure や coarse moduli space の特異点の同定ではない．
\end{proof}

\begin{corollary}[非零 Poisson 構造の射影化]
\label{cor:projective-moduli}
$\mathcal P_{n,d}^{\times}$ を零切断の補集合とする．
Poisson 構造への独立なスカラー作用による商
$[\mathcal P_{n,d}^{\times}/\mathbb G_m]$ から $\mathcal D_{n,d}$ への射は，
representable，滑らかかつ射影的である．
線束 $\mathcal H$ を選べる任意の基底上で，この射は
\[
 \operatorname{Gr}_T(n-1,p_*\mathcal H)\longrightarrow T
\]
に同型である．$\mathcal H$ を基底の線束で捻ってもこの Grassmann 族は変わらない．
\end{corollary}

\begin{proof}
命題~\ref{prop:intrinsic-bivector-bundle} の捻られた Pl\"ucker 錐を射影化すると
$\mathcal L$ の因子が消え，上記の Grassmann 族となる．
式~\eqref{eq:intrinsic-bivector-bundle} のベクトル束は既に $\mathcal D_{n,d}$ 上へ
降下しているため，この族はその射影束の閉部分スキームとして実現され，射影的である．
representable 性と滑らかさも atlas 上の Grassmann 族から従う．
\end{proof}

ここで構成したものは，零 Poisson 構造を含む滑らかな多様体のモジュライスタックである．
特異 del Pezzo 多様体を加えたコンパクト化，その coarse moduli space の具体的な
構成，および相対量子化の降下は，この構成の結論には含めない．

% !TeX root = ../../poisson-del-pezzo.tex
\section{零 Poisson 構造の特異性と標準的な解消}
\label{sec:singularities}

次数2の相対 Jacobi 錐には，非零部分の滑らかさより強い幾何学的構造がある．
本節では，零 Poisson 構造を保つ座標変換に関して標準的な blow-up を構成し，
モジュライスタックに降下させる．以下では $n\geq2$，$m=n+1$ とし，
$C_m=\operatorname{Cone}(\Gr(2,m))\subset\bigwedge^2\C^m$ と書く．
Grassmann 多様体の Pl\"ucker 埋め込みが射影正規で arithmetically
Cohen--Macaulay であることは標準単項式理論の既知の結果として使う
\cite{Svanes1973,DeConciniProcesi1981}．

\begin{lemma}[Pl\"ucker 錐の局所完全交叉性]
\label{lemma:cone-lci}
$C_m$ の頂点 $0$ において
\[
 \operatorname{edim}_0 C_m=\binom m2,\qquad
 \dim_0 C_m=2m-3,\qquad
 \operatorname{codim}_0 C_m=\frac{(m-2)(m-3)}2.
\]
頂点の定義イデアルの最小生成元数は $\binom m4$ である．
従って $m=3$ では $C_m=\mathbb A^3$，$m=4$ では非退化な二次超曲面の錐，
$m\geq5$ では $C_m$ は頂点で局所完全交叉でない．
\end{lemma}
\begin{proof}
錐の座標環を $R=\C[a_{ij}]/I$ とする．$I$ は斉次二次の Pl\"ucker 式
\[
 q_{ijkl}=a_{ij}a_{kl}-a_{ik}a_{jl}+a_{il}a_{jk}
 \qquad (i<j<k<l)
\]
で生成される．これらは相異なる四つの添字集合ごとに，互いに共有しない
二次単項式を含むので一次独立であり，$I/\mathfrak m I$ の基底をなす．
ゆえに最小生成元数は $\binom m4$ である．
$I$ に一次式がないので，接空間次元は $\binom m2$ である．
$\dim\Gr(2,m)=2(m-2)$ から錐の次元は $2m-3$ となる．
$m=3$ の場合は $I=0$，$m=4$ では唯一の Pl\"ucker 式は
六変数の非退化二次形式である．$m\geq5$ の場合には
\[
 \binom m4>\binom m2-(2m-3)
 =\frac{(m-2)(m-3)}2
\]
が成り立つので，正則局所環における完全交叉の生成元数の必要条件に反する．
\end{proof}

\begin{lemma}[頂点の blow-up]
\label{lemma:cone-blowup}
$Y_m=\Gr(2,m)$，$L=\mathcal O_{Y_m}(-1)$ とおく．
$C_m$ の頂点の blow-up は，標準的に
\[
 \operatorname{Bl}_0(C_m)\simeq\operatorname{Tot}_{Y_m}(L)
\]
である．例外因子は零切断 $Y_m$ であり，その法線束は $L$ である．
この同型は $\operatorname{GL}_m$ および錐のスカラー作用に関して同変である．
\end{lemma}
\begin{proof}
$C_m\setminus\{0\}$ の点は，Pl\"ucker 直線 $[v]\in Y_m$ と
その直線の非零ベクトル $v$ の組に一意に対応する．
従って閉部分スキーム
\[
 \{(v,[\ell])\in C_m\times Y_m\mid v\in\ell\}
\]
は $\operatorname{Tot}(L)$ である．一方，頂点のイデアルは一次座標
$a_{ij}$ 全体で生成され，その blow-up は $C_m$ 上のこれらの座標による
有理写像 $C_m\dashrightarrow\mathbb P(\bigwedge^2\C^m)$ のグラフの閉包である．
$C_m\setminus\{0\}$ ではこの写像の像は $Y_m$ に入り，閉包を取っても
$Y_m$ の斉次 Pl\"ucker 方程式が保たれる．上の閉部分スキームは各
Pl\"ucker チャート上で「一つの座標が自由なスカラー，残りはその座標に
Grassmann チャート座標を掛けたもの」という形だから，グラフ閉包に等しい．
すべての構成が線型作用に関して自然である．
\end{proof}

\begin{theorem}[相対零切断の blow-up による解消]
\label{thm:poisson-resolution}
次数2のモジュライの零切断を
$s:\mathcal D_{n,2}\hookrightarrow\mathcal P_{n,2}$ とおく．
$\mathcal R_{n,2}=\operatorname{Bl}_{s(\mathcal D_{n,2})}\mathcal P_{n,2}$
とすると，次が成り立つ．
\begin{enumerate}
\item $\mathcal R_{n,2}$ は有限型の滑らかな Deligne--Mumford スタックであり，
$\mathcal R_{n,2}\to\mathcal P_{n,2}$ は representable，射影的かつ固有で，
零切断の補集合では同型である．$n\geq3$ のときこれは特異性の解消を与える．
\item 例外因子 $\mathcal E_{n,2}$ は $\mathcal D_{n,2}$ 上の滑らかな相対
Grassmann 族
\[
 \mathcal E_{n,2}\simeq
 [\mathcal P_{n,2}^{\times}/\mathbb G_m]
\longrightarrow\mathcal D_{n,2}
\]
に標準的に同型である．相対 Poisson パラメータのベクトル束を
$\mathcal V$，その射影 Grassmann 部分を $\mathcal G\subset\mathbb P(\mathcal V)$
と書けば，解消は $\mathcal G$ 上の tautological line bundle
$\mathcal O_{\mathcal G}(-1)$ の全空間に一致する．
例外因子の法線束は $\mathcal O_{\mathcal G}(-1)$ である．
\item 任意の $S\to\mathcal D_{n,2}$ に対して，この解消は対応する
相対 Pl\"ucker 錐の零切断の blow-up と可換である．
\end{enumerate}
\end{theorem}
\begin{proof}
命題~\ref{prop:intrinsic-bivector-bundle} と
系~\ref{cor:projective-moduli} により，$\mathcal V$ と $\mathcal G$ は
偏極線束を基底から捻る操作に対して不変であり，$\mathcal D_{n,2}$ 上に降下する．
$\mathcal D_{n,2}$ の標準方程式 atlas $U_2$ 上では，
$\mathcal P_{n,2}$ は $U_2\times C_m$，零切断は $U_2\times\{0\}$ である．
blow-up は平坦基底変換と可換なので，補題~\ref{lemma:cone-blowup} から
\[
 \mathcal R_{n,2}\times_{\mathcal D_{n,2}}U_2
 \simeq U_2\times\operatorname{Tot}_{Y_m}\mathcal O_{Y_m}(-1)
\]
を得る．同型は式~\eqref{eq:moduli-action-two} の $G_2$ の作用と同変である．
従って商スタックへ降下し，滑らかさ，例外因子と法線束の記述を得る．
blow-up の射影性・固有性と零切断外での同型性も atlas 上で確認して降下する．
記載した基底変換は各 atlas 上での自明な積表示と blow-up の構成から従う．
\end{proof}

\begin{proposition}[Pl\"ucker 錐の標準特異性]
\label{prop:plucker-terminal}
$m\geq4$ のとき $C_m$ は頂点のみで特異であり，その特異点は
Gorenstein terminal（従って rational）である．
補題~\ref{lemma:cone-blowup} の解消 $\rho:\widetilde C_m\to C_m$ の例外因子を
$E$ とすると
\[
 K_{\widetilde C_m}=\rho^*K_{C_m}+(m-1)E
\]
という discrepancy 等式が成立する．特に $m=n+1$ では例外因子の
discrepancy は $n$ である．これは Poisson モジュライの滑らかな
標準方程式 atlas に関する主張であって，coarse moduli space の
商特異点についての主張ではない．
\end{proposition}
\begin{proof}
Grassmann 多様体の普遍部分束 $\mathcal S$ と商束 $\mathcal Q$ から
$T_{Y_m}=\mathcal S^*\otimes\mathcal Q$，従って
$\omega_{Y_m}\simeq\mathcal O_{Y_m}(-m)$ である．
Pl\"ucker 座標環 $R$ は射影正規かつ Cohen--Macaulay であり，
その次数付き canonical module は
\[
 \omega_R\simeq
 \bigoplus_{k\in\mathbb Z}H^0(Y_m,\omega_{Y_m}(k))
 \simeq R(-m)
\]
である．ここで最後の等式では $\omega_{Y_m}=\mathcal O(-m)$ と
Pl\"ucker 線束の切断環の記述を用いた．従って $C_m$ は Gorenstein であり，
$K_{C_m}$ は（非斉次な意味で）自明である．

$\widetilde C_m=\operatorname{Tot}_{Y_m}(L)$ に対して
$\omega_{\widetilde C_m}\simeq
p^*(\omega_{Y_m}\otimes L^{-1})=p^*\mathcal O_{Y_m}(-(m-1))$ である．
零切断 $E$ の法線束が $L=\mathcal O(-1)$ なので
$\mathcal O_{\widetilde C_m}(E)=p^*\mathcal O(-1)$ であり，
discrepancy は $m-1$ である．頂点以外の錐は滑らかだから，
他の例外因子についても滑らかな解消の上でさらに吹き上げれば
discrepancy が正になる．従って terminal である．
標数0で terminal は klt，klt は rational である
\cite[Chapter 5]{KollarMori1998}．
\end{proof}

\begin{corollary}[相対モジュライの特異性の次元別分類]
\label{cor:moduli-singularity-types}
次数2の $\mathcal P_{n,2}\to\mathcal D_{n,2}$ は，
$n=2$ では滑らか，$n=3$ では零切断に沿い相対的に通常の
五次元二次錐特異点を持つ hypersurface 型，
$n\geq4$ では零切断に沿い非局所完全交叉である．
$n\geq3$ の零点における相対 Zariski 接空間の次元超過は
\[
 \binom{n+1}{2}-(2n-1)=\frac{(n-1)(n-2)}2
\]
である．どの $n\geq3$ でも滑らかな atlas の特異性は terminal であり，
その零切断 blow-up は定理~\ref{thm:poisson-resolution} の滑らかな解消となる．
\end{corollary}
\begin{proof}
相対分類，補題~\ref{lemma:cone-lci}，
定理~\ref{thm:poisson-resolution}，命題~\ref{prop:plucker-terminal} を合わせる．
\end{proof}

\begin{remark}[既知の Grassmann 幾何と本稿の寄与]
錐の blow-up と Gorenstein 性，標準単項式理論に基づく局所環の性質は
既知の代数幾何である．ここでの寄与は，次数2の相対 Jacobi スキームとの
スキーム論的同定を経由し，零 Poisson 切断とその解消を
$G_2$-同変に具体化して，Poisson del Pezzo のモジュライスタックに
降下させた点にある．この降下による新規性の優先権は未確認である．
\end{remark}


\section{零 Poisson 構造の無限小変形}
\label{sec:infinitesimal}
スキームとしての分類は、複素点の分類だけでは見えない差を明示する。
次数1では全ての変形方向が可換なベクトル空間をなし、任意の階数の Artin 基底上でも
Jacobi 方程式による制約はない。次数2では零点で二次の障害が現れる。

\begin{corollary}
\label{cor:second-order}
次数2の座標同定 $V_X\simeq\bigwedge^2\C^{n+1}$ を固定する。
任意の $Q\in V_X$ に対し $\varepsilon Q$ は
$\C[\varepsilon]/(\varepsilon^2)$ 上の零 Poisson 構造の変形である。
これが $\C[\varepsilon]/(\varepsilon^3)$ 上へ持ち上がる必要十分条件は
$Q\wedge Q=0$ である。この条件の下では $tQ$ により全次数へ持ち上がる。
\end{corollary}
\begin{proof}
Pl\"ucker 方程式は斉次二次なので、双対数上では自動的に消える。
任意の二次の持ち上げは $\varepsilon Q+\varepsilon^2R$ と書け、
その外積二乗は $\varepsilon^2Q\wedge Q\pmod{\varepsilon^3}$ である。
従って $R$ の選び方では障害を除けない。
$Q\wedge Q=0$ なら $tQ$ が Jacobi 方程式を恒等的に満たす。
\end{proof}

例えば $n\geq3$ で
$Q=\partial_{x_0}\wedge\partial_{x_1}+\partial_{x_2}\wedge\partial_{x_3}$
とすると、これは一次変形を与えるが二次へ持ち上がらない。
$Q\wedge Q=2\partial_{x_0}\wedge\partial_{x_1}\wedge
\partial_{x_2}\wedge\partial_{x_3}\neq0$ だからである。
この例は、$\dim T_0\Pois_{\mathrm{sch}}(X)=\binom{n+1}{2}$ が
実際の次元 $2n-1$ より大きいことを具体化する。
曲面の場合には $\bigwedge^4\C^3=0$ なのでこの障害はない。

% !TeX root = ../poisson-del-pezzo.tex
\section{次数付き形式量子化と中心的超曲面}
\label{sec:quantization}

この節では、第~\ref{sec:classification}節で分類した各 Poisson 構造について、標準重み付き多項式環を
量子化し、定義式を中心元に持ち上げる。
外部入力は Kontsevich の affine formality とその twisted 接写像である
\cite[\S\S 4.6.2, 6, 8]{Kontsevich}。
特異な超曲面 cone に直接 formality を適用するのではなく、滑らかな affine
空間 $\operatorname{Spec}S$ に適用する。

\subsection{重みと中心的線型座標を保つ積}

$R=\mathbb C[[\hbar]]$ とおき、内部次数を $\deg\hbar=0$ とする。
$S$ は標準 ambient の重み付き多項式環、$F\in S_m$ は定義式であり、
$m=6$（次数 $1$）または $m=4$（次数 $2$）である。
既述の $F$-Casimir 代表を $\widetilde\pi$ と書く。

\begin{theorem}[中心的超曲面による形式量子化]
\label{thm:formal-quantization}
各 Poisson 構造 $\pi$ に対し、次数付き $R$-代数 $\mathcal S_\hbar$ と
斉次中心元 $\Omega_\hbar\in Z(\mathcal S_\hbar)_m$ で、次を満たすものが存在する。
\begin{enumerate}
\item 各内部次数 $\ell$ の成分 $(\mathcal S_\hbar)_\ell$ は有限自由
      $R$-加群であり、$\mathcal S_\hbar/\hbar\mathcal S_\hbar\simeq S$ である。
\item $\Omega_\hbar\bmod\hbar=F$ であり、$\Omega_\hbar$ は両側非零因子である。
\item $\mathcal R_\hbar:=\mathcal S_\hbar/(\Omega_\hbar)$ も各内部次数で有限自由で、
      $\mathcal R_\hbar/\hbar\mathcal R_\hbar\simeq S/(F)$ である。
      この商の半古典括弧は $\pi$ を誘導する。
\item 非零 $\pi$ に対し、既述の $n-1$ 個の中心的線型座標 $s_i$ は
      補正なしに中心的に保てる。さらに積は $\mathbb C[s_1,\ldots,s_{n-1}]$-双線型に選べる。
\end{enumerate}
これは存在定理であり、量子化の同型類や gauge 同値類の一意性を主張しない。
\end{theorem}

\begin{proof}
Affine formality の座標公式によって
\begin{equation}
 a\star b=ab+\frac{\hbar}{2}\{a,b\}+O(\hbar^2),
 \qquad a\star b-b\star a=\hbar\{a,b\}+O(\hbar^2)
 \label{eq:star-first-order}
\end{equation}
という結合積を取る。各 $\hbar$ 係数は有限個のグラフから得られる
多項式係数の有限階双微分作用素である。
複素多項式に対する同じ公式は、実係数の普遍的グラフ定数を複素数体に移して用いる。
$\widetilde\pi$ の重みが $0$ なので、各微分の重みとその添字に対応する
多重ベクトル係数の重みが相殺し、全項が内部次数を保つ。
従って
\begin{equation}
 \mathcal S_\hbar=\bigoplus_{\ell\geq0}S_\ell[[\hbar]]
 \label{eq:graded-formal-ambient}
\end{equation}
は $\star$ で閉じる。ここでは内部次数について直和を取り、各成分だけを
$\hbar$ 進完備化する。$\hbar$ の次数とともに内部次数が無限に増大する元を含む
$S[[\hbar]]$ 全体とは区別する。

非零構造については、活性な三変数を $u,v,w$、残りの線型座標を $s_i$ と取ると、
$\widetilde\pi$ に $\partial_{s_i}$ 成分はない。
従ってグラフ公式の各矢印は活性変数の方向にのみ微分する。
正規化した unital な公式では
\begin{equation}
 s_i\star a=s_i a=a\star s_i,
 \qquad
 (g(s)a)\star b=g(s)(a\star b)=a\star(g(s)b)
 \label{eq:exact-central-parameters}
\end{equation}
となる。ここで係数として現れる $s_i$ を微分する矢印もないことを用いた。

次に $\gamma=\hbar\widetilde\pi$ とし、formality 写像を $\mathcal U$ と書く。
その twist の関数成分は
\begin{equation}
 Q_\gamma(f)=\sum_{r\geq0}\frac{1}{r!}
   \mathcal U_{r+1}(f,\underbrace{\gamma,\ldots,\gamma}_{r\text{ 個}})
 =f+O(\hbar)
 \label{eq:twisted-central-map}
\end{equation}
である。Twisted chain map の性質により $[\widetilde\pi,f]=0$ ならば
$Q_\gamma(f)$ は Hochschild の $0$-cocycle、すなわち $\star$ 積の中心元になる。
この次数では coboundary を法とすることによる曖昧さはなく、実際の中心元が得られる。
従って
\begin{equation}
 \Omega_\hbar=Q_\gamma(F)\in Z(\mathcal S_\hbar)_m,
 \qquad \Omega_\hbar=F+O(\hbar)
 \label{eq:central-equation}
\end{equation}
を取れる。重みの保存により内部次数は $m$ のままである。
一般に $F$ 自体が量子化後も中心的とは限らず、高次補正は省略できない。
非零因子性と商の平坦性は次の補題から従う。
零 Poisson 構造には可換積と $\Omega_\hbar=F$ を用いる。
\end{proof}

\begin{lemma}[中心元を持ち上げた商の平坦性]
\label{prop:central-flatness}
$A=\bigoplus_{\ell\geq0}A_\ell$ は各 $A_\ell$ が有限自由 $R$-加群である
次数付き $R$-代数とする。斉次中心元 $\Omega\in A_m$ の還元
$f\in A/\hbar A$ が非零因子ならば、$\Omega$ は両側非零因子であり、
$A/(\Omega)$ は各内部次数で有限自由 $R$-加群である。
\end{lemma}

\begin{proof}
$\Omega a=0$ を $\hbar$ で還元すると $fa_0=0$ なので $a\in\hbar A$ となる。
これを反復し、各内部次数の $\hbar$ 進分離性を用いると $a=0$ である。
中心性により右側も同様である。
また $\hbar a=\Omega b$ なら、還元して $fb_0=0$ だから $b=\hbar b_1$ と書ける。
$A$ の $\hbar$-torsion freeness により $a=\Omega b_1$ となるので、
$A/(\Omega)$ に $\hbar$-torsion はない。
各内部次数の商は有限生成 $R$-加群であり、DVR 上の有限生成 torsion free
加群は自由である。還元と商が交換することは表示から従う。
\end{proof}

\subsection{Hilbert 級数、有限段階の射影化、中心的切断}

$R$ 上の rank を用いた Hilbert 級数は
\begin{equation}
 H_{\mathcal R_\hbar}(t)=
 \begin{cases}
 \displaystyle\frac{1-t^6}{(1-t)^n(1-t^2)(1-t^3)}&d=1,\\[6pt]
 \displaystyle\frac{1-t^4}{(1-t)^{n+1}(1-t^2)}&d=2.
 \end{cases}
 \label{eq:quantized-hilbert}
\end{equation}
実際、$\Omega_\hbar$ による掛け算の完全列
$0\to\mathcal S_\hbar(-m)\to\mathcal S_\hbar\to\mathcal R_\hbar\to0$
を各内部次数で取ればよい。

各 $r\geq1$ に対し、有限段階の射影的対象を
\begin{equation}
 \operatorname{qgr}(\mathcal R_\hbar/\hbar^r\mathcal R_\hbar)
 \label{eq:formal-qgr-tower}
\end{equation}
で表す。ここで $\operatorname{qgr}$ は有限生成次数付き右加群の圏を
右有界加群の Serre 部分圏で割ったものである。
各有限段階は左右 noetherian である。実際、有限の $\hbar$ 進 filtration の
随伴次数環は $(S/(F))\otimes_\mathbb C\mathbb C[\hbar]/(\hbar^r)$ と同型であり、
noetherian 性を有限 filtration に沿って持ち上げられる。
$r=1$ では、$H$ の section ring と Serre の対応から
$\operatorname{qgr}(S/(F))\simeq\operatorname{coh}(X)$ である。
式\eqref{eq:formal-qgr-tower}は還元に関する整合的な系を与えるが、
この記述だけで任意の基底上の Abel 圏変形の普遍性を主張するものではない。

式\eqref{eq:exact-central-parameters}により、$s_i$ の斉次線型関係による切断を
量子化後にも取れる。関係を課した可換極限で $F$ が非零多項式のままであるなら、
補題\ref{prop:central-flatness}をその多項式環に適用し、平坦な量子化を得る。
特に、射影的な共通 level に対応する線型関係を課すことで、
既述の二次元 leaf の閉包に対応する曲面切断を量子化できる。
この操作は中心的線型関係を先に課した後の多項式環に対する議論であり、
一般の非平坦な基底変換との交換を無条件に述べているわけではない。

\section{曲面の cyclic potential と Rees コンパクト化}
\label{sec:surface-potential}

以下では $n=2$ とし、非零構造に対応する中心的線型座標を $s$ とする。
活性変数 $u,v,w$ の重みを
\[
 (a,b,c)=
 \begin{cases}(1,2,3)&d=1,\\(1,1,2)&d=2,\end{cases}
 \qquad m=a+b+c
\]
と書く。$s\ne0$ 上の座標を $u/s^a,v/s^b,w/s^c$ とし、同じ文字で略記すると
\begin{equation}
 U=X\setminus(s=0)=\operatorname{Spec}\bigl(\mathbb C[u,v,w]/(\phi)\bigr),
 \qquad \phi(u,v,w)=F(1,u,v,w),\qquad \deg\phi\leq m.
 \label{eq:affine-surface}
\end{equation}
括弧の向きを
$\{u,v\}=\phi_w$, $\{v,w\}=\phi_u$, $\{w,u\}=\phi_v$ に固定する。

\subsection{Potential 表示と中心元}

自由結合代数 $T=\mathbb C\langle u,v,w\rangle$ の cyclic quotient を
$T_{\mathrm{cyc}}=T/[T,T]$ とする。語に対する cyclic derivative の規約は
\begin{equation}
 \partial_u^{\mathrm{cyc}}[a_1\cdots a_r]
 =\sum_{i:a_i=u}a_{i+1}\cdots a_r a_1\cdots a_{i-1}
 \label{eq:cyclic-derivative}
\end{equation}
である。$\phi$ の重み $m$ 以下の cyclic lift $\widehat\phi$ を選び、
\begin{equation}
 \Phi_\hbar=[uvw-vuw]-\hbar\widehat\phi,
 \qquad
 \mathcal A_\hbar=
 \frac{\mathbb C\langle u,v,w\rangle[[\hbar]]}
 {\overline{(\partial_u^{\mathrm{cyc}}\Phi_\hbar,
             \partial_v^{\mathrm{cyc}}\Phi_\hbar,
             \partial_w^{\mathrm{cyc}}\Phi_\hbar)}}
 \label{eq:potential-algebra}
\end{equation}
と定める。上線は $\hbar$ 進閉包を表し、語長については完備化しない。
直接微分すると、関係式は
\begin{equation}
 [u,v]=\hbar\partial_w^{\mathrm{cyc}}\widehat\phi,
 \quad[v,w]=\hbar\partial_u^{\mathrm{cyc}}\widehat\phi,
 \quad[w,u]=\hbar\partial_v^{\mathrm{cyc}}\widehat\phi
 \label{eq:potential-relations}
\end{equation}
となり、半古典括弧の符号が上の規約に一致する。

\begin{proposition}[Etingof--Ginzburg の結果の適用]
\label{prop:eg-potential}
$\mathcal A_\hbar$ は $\mathbb C[u,v,w]$ の平坦な形式変形である。
重み filtration が $m$ 以下の中心元 $\Psi_\hbar$ で
$\Psi_\hbar\bmod\hbar=\phi$ を満たすものが存在する。
従って $\mathcal B_\hbar=\mathcal A_\hbar/(\Psi_\hbar)$ は
$\mathbb C[U]$ の平坦な形式量子化である。
\end{proposition}

\begin{proof}
平坦性には \cite[Lemma 7.4.2]{EG2010} を適用する。
必要なのは potential の定数項が $uvw-vuw$ で、残りが $\hbar$ で割れることである。
中心元には \cite[Proposition 8.1.1(i)]{EG2010} を適用する。
その仮定は三つの重みが正であり、potential の重みがその和 $m$ 以下であることなので、
本件で満たされる。同命題の中心元の還元は potential の $\hbar$ 一次項の
abelianization である。本稿ではその一次項が $-\widehat\phi$ なので、
得られた中心元の負を $\Psi_\hbar$ と取る。
商の平坦性は、$\phi$ が多項式環の非零因子であることと
補題\ref{prop:central-flatness}の $\hbar$ 進分離性による議論から従う。
Filtered な各重み成分でも、同文献の Rees 構成によりこの議論を行える。
\end{proof}

\begin{remark}
Proposition 8.1.1(i) 自体には $\phi$ の最高重み部分の孤立特異点という仮定はない。
同命題の追加の中心・Hochschild コホモロジーの結論 (ii), (iii) には
斉次性と孤立特異点の仮定がある。
従って上の中心元の存在は反標準境界が特異な場合にも適用できるが、
その場合に中心全体が一元で生成されるとは結論しない。
また、この potential による量子化と定理\ref{thm:formal-quantization}で選んだ
量子化が同じ gauge 同値類に属することも、ここでは主張しない。
\end{remark}

\subsection{Rees 代数と二種類の中心元}

重み filtration を $\mathcal F_\ell\mathcal A_\hbar$ と書く。
各 filtration 成分を $\hbar$ 進完備化した上で、内部次数について直和を取り、
\begin{equation}
 A_\hbar=\bigoplus_{\ell\geq0}
       (\mathcal F_\ell\mathcal A_\hbar)s^\ell,
 \qquad\deg s=1
 \label{eq:rees-algebra}
\end{equation}
とする。Rees 生成元 $us^a,vs^b,ws^c$ を再び $u,v,w$ と書く。
$\widehat\phi$ の重み $m$ への斉次化を $\widehat F$ とすれば、この代数は
$s$ を中心とする $\mathbb C[s]$ 上の相対的三変数 potential
\[
 \Phi^{\mathrm{hom}}_\hbar=[uvw-vuw]-\hbar\widehat F(s,u,v,w)
\]
により表示される。課すのは $u,v,w$ の cyclic derivative だけであり、
$s$ に関する cyclic derivative は課さない。
\cite[Proposition 8.1.1(i) の証明、pp.~1407--1408]{EG2010}の相対的 Rees 構成により、
$A_\hbar$ は $\mathbb C[s,u,v,w]$ の次数付き平坦変形である。

$\Psi_\hbar$ を次数 $m$ に斉次化した中心元を $\Omega_\hbar$ とすると
\begin{equation}
 B_\hbar=A_\hbar/(\Omega_\hbar),
 \qquad B_\hbar/\hbar B_\hbar\simeq\mathbb C[s,u,v,w]/(F)
 \label{eq:rees-hypersurface}
\end{equation}
を得る。ここには役割の異なる二種類の中心元がある。
$s$ は反標準境界を指定し、$s=1$ はその affine 補集合を与える。
$\Omega_\hbar$ は非可換 ambient 内の超曲面そのものを定義する。

$s$ は Rees 代数 $A_\hbar$ 上の非零因子であり、
$A_\hbar/(s-1)$ の $\hbar$ 進完備化は $\mathcal A_\hbar$ である。
また、$\Omega_\hbar\bmod s$ のさらに $\hbar$ による還元は非零多項式
$F(0,u,v,w)$ である。従って $A_\hbar/(s)$ で $\Omega_\hbar\bmod s$ は
非零因子となる。$sa=\Omega_\hbar b$ を $s$ で還元すると $b\in sA_\hbar$ なので、
$s$ を消去して $a\in\Omega_\hbar A_\hbar$ を得る。
従って $s$ は $B_\hbar$ 上でも非零因子である。

\subsection{具体的な関係式}

次数 $1$ では
\[
 F=w^2-v^3-A_4(s,u)v-A_6(s,u),
 \qquad (\deg s,\deg u,\deg v,\deg w)=(1,1,2,3)
\]
と書く。$A_4(s,u)=\sum_{k=0}^4a_k s^{4-k}u^k$ とし、$A_4(s,u)v$ を
この順序の cyclic 語で持ち上げると
\begin{align}
 [u,v]&=2\hbar w,\label{eq:degree1-uv}\\
 [w,u]&=-\hbar\bigl(3v^2+A_4(s,u)\bigr),\label{eq:degree1-wu}\\
 [v,w]&=-\hbar\left(
    \sum_{k=1}^4 a_k s^{4-k}\sum_{j=0}^{k-1}u^{k-1-j}vu^j
    +\frac{\partial A_6}{\partial u}(s,u)\right),\label{eq:degree1-vw}\\
 [s,u]&=[s,v]=[s,w]=0.\nonumber
\end{align}
実際、$\partial_u^{\mathrm{cyc}}(u^kv)
=\sum_{j=0}^{k-1}u^{k-1-j}vu^j$ である。
非可換代数内ではこの和を $ku^{k-1}v$ に置き換えられない。
超曲面にはさらに次数 $6$ の中心元 $\Omega_\hbar=F+O(\hbar)$ を課す。

次数 $2$ では
\[
 F=w^2-f_4(s,u,v),
 \qquad(\deg s,\deg u,\deg v,\deg w)=(1,1,1,2)
\]
と書く。$\widehat f_4$ を重み $4$ の cyclic lift とすれば
\begin{equation}
 \begin{gathered}
 [u,v]=2\hbar w,\qquad
 [v,w]=-\hbar\partial_u^{\mathrm{cyc}}\widehat f_4,\qquad
 [w,u]=-\hbar\partial_v^{\mathrm{cyc}}\widehat f_4,\\
 [s,u]=[s,v]=[s,w]=0
 \end{gathered}
 \label{eq:degree2-potential-relations}
\end{equation}
となる。超曲面の方程式は次数 $4$ の中心元 $\Omega_\hbar=F+O(\hbar)$ である。
これらの表示は中心元の高次補正そのものの閉じた公式を与えるものではない。

\section{AS 正則性、型の呼称、分類の限界}
\label{sec:quantization-limits}

\subsection{中心拡大に関する条件付き結論}

体 $k$ 上の connected な正次数付き代数に対し、AS 正則性には有限大域次元、
多項式増大度、左右の AS--Gorenstein 条件を要求する。
$R=\mathbb C[[\hbar]]$ 上の形式平坦性とは別の性質である。

\begin{proposition}[中心拡大]
\label{prop:central-as-extension}
$A$ を connected、局所有限、左右 noetherian な次数付き $k$-代数とし、
$s\in A_1$ を中心的非零因子とする。
$C=A/(s)$ が大域次元 $3$ の AS 正則代数ならば、$A$ は大域次元 $4$ の
AS 正則代数である。さらに $\Omega\in Z(A)_m$ が非零因子ならば
$A/(\Omega)$ はこの AS 正則 ambient の中心的超曲面環である。
\end{proposition}

\begin{proof}
完全列 $0\to A(-1)\xrightarrow{s}A\to C\to0$ により、
$C$ は左右の $A$-加群として射影次元 $1$ を持つ。
$k$ の長さ $3$ の $C$-自由分解の各項を $A$ 上で長さ $1$ の分解に置き換えれば、
$\operatorname{pd}_A k\leq4$ である。
ここから全加群の大域次元の上界を得る議論も記しておく。
$A^e=A\otimes_k A^{\mathrm{op}}$ 上の $A$ の graded minimal 自由分解を取る。
二側 bar 分解に augmentation 加群 $k$ をテンソルすると、通常の reduced bar
複体を得るので
$\operatorname{Tor}^{A^e}_i(k,A)\simeq\operatorname{Tor}^A_i(k,k)$ である。
右辺は $i>4$ で消え、minimal 性と graded Nakayama により上の
$A^e$-自由分解の第 $i$ 項も消える。
その分解を右 $A$-加群の分解と見れば、終項 $A$ が射影的なので分裂完全である。
従って任意の左 $A$-加群との右側テンソルは完全性を保ち、長さ $4$ 以下の
左自由分解を与える。左右を入れ替えた議論も同様であり、
$\operatorname{gldim}A\leq4$ を得る。
一方、同じ完全列に $\operatorname{Hom}_A(-,A)$ を適用すると
$\operatorname{Ext}_A^q(C,A)=0$（$q\ne1$）、
$\operatorname{Ext}_A^1(C,A)\simeq C(1)$ である。
環の変更の Ext スペクトル系列は一行に退化し、
\[
 \operatorname{Ext}_A^{p+1}(k,A)
 \simeq \operatorname{Ext}_C^p(k,C)(1)
\]
を与える。$C$ の AS--Gorenstein 条件から、この群は次数 $4$ だけで一次元となる。
左加群側も同様なので、左右の AS--Gorenstein 条件と大域次元 $4$ を得る。
最後に $H_A(t)=H_C(t)/(1-t)$ から多項式増大度も従う。
\end{proof}

この命題を特定の複素代数または $\mathbb C((\hbar))$ 上の代数に適用するには、
その係数体上で $A/(s)$ の AS 正則性と $A$ の noetherian 性などの仮定を確認する。
形式平坦性や Hilbert 級数の一致だけで、全ての複素パラメータの特殊化の
AS 正則性が得られるわけではない。
上の条件が満たされ、可換模型と同じ Hilbert 級数を持つ場合の ambient は
\[
 H_A(t)=
 \begin{cases}
 ((1-t)^2(1-t^2)(1-t^3))^{-1}&d=1,\\
 ((1-t)^3(1-t^2))^{-1}&d=2,
 \end{cases}
\]
であり、曲面をそれぞれ $\operatorname{qgr}(A/(\Omega_6))$、
$\operatorname{qgr}(A/(\Omega_4))$ と表す。
AS 正則性は ambient の $A$ に対する主張であり、超曲面環 $A/(\Omega)$ の
AS 正則性を述べているのではない。

境界側 $C=A/(s)$ で $\hbar$ が可逆ならば、両次数で
$w=[u,v]/(2\hbar)$ を消去できる。
次数 $2$ は重み $1$ の二生成元と二個の三次関係式の表示、
次数 $1$ は重み $(1,2)$ の二生成元と重み $4,5$ の関係式の表示になる。
これは cubic 型および Stephenson 型の正則代数との接点を示すが、
生成元・関係式の重みだけによって正則性を判定することはできない。

\subsection{\texorpdfstring{$E_7,E_8$}{E7, E8} 型と非形式的な族}

活性変数の重み $(1,1,2)$（次数 $2$）と $(1,2,3)$（次数 $1$）は、
\cite[\S\S2.3, 3.4]{EG2010}の $E_7,E_8$ 型に対応する。
反標準境界 $(s=0)$ が滑らかな楕円曲線なら、
$F(0,u,v,w)=0$ はそれぞれ simple elliptic 型 $\widetilde E_7,\widetilde E_8$
の重み付き cone である。滑らかな曲面 $X$ 自身にその型の特異点があるという意味ではない。
境界が特異なら、その cone に孤立特異点の結果をそのまま適用できない。

Etingof--Ginzburg は適切な正規形の potential
\[
 [uvw-qvuw]-\lambda\bigl(P(u)+Q(v)+R(w)\bigr)
\]
を扱う。最高次数の組は $E_7$ 型で $(4,4,2)$、$E_8$ 型で $(6,3,2)$ である。
\cite[Theorem 3.3.2]{EG2010} は最高重みの斉次 potential が三次元 CY であることを
仮定して、低重み項を加えた filtered 代数について結論する。
同論文の Theorems 3.4.4--3.4.5 による正規形、semiuniversal な族、中心の記述には
genericity がある。これらを全パラメータでの正則性に読み替えない。

形式級数を任意の非零複素数 $\hbar$ に代入することは、収束または代数化を
別に示さない限り定義されない。
Poisson 構造は変形の一次反対称部分を指定するのであり、
高次の potential、中心元の高次補正、座標変換および gauge 同値まで固定するわけではない。
従って、本稿が与える全 Poisson 構造の中心的超曲面量子化の存在から、
$\operatorname{coh}(X)$ の全ての非可換変形の網羅性や intrinsic なモジュライの
同定は従わない。この逆向きの認識、相対的下降、代数化、特殊化の正則性は
それぞれ独立した研究課題である。


\section{今後の数学的課題}
\label{sec:open}
本稿では固定した $X$ の Poisson スキームから動く $X$ の相対分類と
モジュライスタックまでを構成した。
coarse moduli space の具体的構成、GIT 安定性、特異な del Pezzo 多様体を含む
コンパクト化は、ここでの商スタックの構成とは別の課題である。
また全 Poisson パラメータ空間上の量子化の選択とその降下、代数化、
全ての非可換変形の網羅性は未解決である。
曲面の非形式的な量子化では、中心元の補正、Rees algebra、
genericity と AS 正則性の仮定を独立に検証する必要がある。
これらを本稿で証明済みの結果と混同しない。

\begingroup
\small
\bibliographystyle{alpha}
\bibliography{references}
\endgroup
\end{document}
